\documentclass[11pt,letterpaper]{article}

\usepackage[utf8]{inputenc}
\usepackage[T1]{fontenc}
\usepackage[spanish,es-nodecimaldot]{babel}
\usepackage{lmodern}
\usepackage{microtype}
\usepackage{geometry}
\geometry{margin=2.45cm}
\usepackage{setspace}
\setstretch{1.08}
\usepackage{graphicx}
\usepackage{booktabs}
\usepackage{tabularx}
\usepackage{array}
\usepackage{longtable}
\usepackage{multirow}
\usepackage{siunitx}
\sisetup{detect-all=true,round-mode=places,round-precision=3}
\usepackage{amsmath,amssymb,mathtools}
\usepackage{caption}
\usepackage{subcaption}
\usepackage{float}
\usepackage{placeins}
\usepackage{xcolor}
\usepackage[most]{tcolorbox}
\usepackage{enumitem}
\usepackage{hyperref}
\usepackage[numbers,sort&compress]{natbib}
\usepackage{fancyhdr}
\usepackage{titlesec}
\usepackage{pdflscape}

\definecolor{cmatblue}{RGB}{31,78,121}
\definecolor{cmatlight}{RGB}{238,245,250}
\definecolor{cmatgray}{RGB}{245,245,245}
\definecolor{cmatgreen}{RGB}{235,247,238}
\definecolor{cmatorange}{RGB}{253,243,229}

\hypersetup{
  colorlinks=true,
  linkcolor=cmatblue,
  citecolor=cmatblue,
  urlcolor=cmatblue,
  pdftitle={Del uso incentivado a la búsqueda persistente de apoyo: resultados preliminares CMAT},
  pdfauthor={Centro de Aprendizaje de Matemáticas, UDLAP}
}

\pagestyle{fancy}
\fancyhf{}
\lhead{CMAT -- UDLAP}
\rhead{Borrador de resultados}
\cfoot{\thepage}
\setlength{\headheight}{14pt}

\titleformat{\section}{\Large\bfseries\color{cmatblue}}{\thesection}{0.7em}{}
\titleformat{\subsection}{\large\bfseries}{\thesubsection}{0.7em}{}
\titleformat{\subsubsection}{\normalsize\bfseries}{\thesubsubsection}{0.7em}{}

\newtcolorbox{questionbox}[1][]{
  enhanced,
  breakable,
  colback=cmatlight,
  colframe=cmatblue,
  boxrule=0.6pt,
  arc=1.5mm,
  title=Pregunta,
  fonttitle=\bfseries,
  #1
}
\newtcolorbox{answerbox}[1][]{
  enhanced,
  breakable,
  colback=cmatgreen,
  colframe=green!45!black,
  boxrule=0.6pt,
  arc=1.5mm,
  title=Respuesta empírica,
  fonttitle=\bfseries,
  #1
}
\newtcolorbox{interpretbox}[1][]{
  enhanced,
  breakable,
  colback=cmatorange,
  colframe=orange!60!black,
  boxrule=0.6pt,
  arc=1.5mm,
  title=Qué significa y qué podría estar pasando,
  fonttitle=\bfseries,
  #1
}
\newtcolorbox{cautionbox}[1][]{
  enhanced,
  breakable,
  colback=cmatgray,
  colframe=black!45,
  boxrule=0.5pt,
  arc=1.5mm,
  title=Precaución inferencial,
  fonttitle=\bfseries,
  #1
}

\newcommand{\CI}{IC\,95\%}
\newcommand{\zscore}{Z}
\newcommand{\ppa}{PPA}
\newcommand{\CMAT}{CMAT}
\newcommand{\MU}{Matemáticas Universitarias}

\begin{document}

\begin{titlepage}
\centering
\vspace*{1.5cm}
{\Large Universidad de las Américas Puebla\par}
\vspace{0.35cm}
{\Large Centro de Aprendizaje de Matemáticas (CMAT)\par}
\vspace{1.8cm}
{\Huge\bfseries Del uso incentivado a la búsqueda persistente de apoyo\par}
\vspace{0.5cm}
{\LARGE Utilización longitudinal del CMAT y desempeño académico\par}
\vspace{0.8cm}
{\large Borrador de resultados y discusión del estudio orientado a publicación\par}
\vfill
\begin{tcolorbox}[colback=cmatlight,colframe=cmatblue,width=0.88\textwidth,boxrule=0.7pt]
\centering
\textbf{Outcome principal:} desempeño académico continuo estandarizado por profesor y periodo.\\[0.2cm]
\textbf{Cohorte principal:} primer intento de Matemáticas Universitarias.\\[0.2cm]
\textbf{Exposición:} cualquier visita al CMAT durante el mismo semestre del primer intento.\\[0.2cm]
\textbf{Regla institucional:} el punto PPA se obtiene al alcanzar tres visitas al CMAT.
\end{tcolorbox}
\vfill
{\large Agosto de 2026\par}
\end{titlepage}

\pagenumbering{roman}
\begin{abstract}
Este trabajo analiza la relación entre la utilización del Centro de Aprendizaje de Matemáticas (CMAT), el desempeño académico continuo y la persistencia posterior de la búsqueda de apoyo. A diferencia del reporte exploratorio previo, la exposición se reconstruye temporalmente: para cada estudiante se cuentan únicamente las visitas al CMAT realizadas durante el mismo periodo de su primer intento de Matemáticas Universitarias. La cohorte primaria contiene 6,627 estudiantes en periodos con cobertura de asesorías. El 18.6\% registró al menos una visita, el 7.1\% alcanzó el umbral institucional de tres visitas requerido para obtener un punto PPA y el 4.6\% continuó más allá de ese umbral.

El outcome principal es la calificación final estandarizada dentro de profesor por periodo. El desempeño medio aumenta de $-0.060$ desviaciones estándar entre quienes no visitan el CMAT a $0.218$ entre quienes registran 1--2 visitas, $0.325$ entre quienes registran exactamente 3 y $0.342$ entre quienes registran 4 o más. El contraste $V>3$ frente a $V\leq3$ produce una diferencia no ajustada de $0.359$ desviaciones estándar (\CI\ $[0.277,0.439]$); con efectos fijos profesor-periodo y ajuste por carrera, la asociación estimada es $0.375$ (\CI\ $[0.257,0.493]$). Sin embargo, la comparación por cuatro grupos muestra una posible meseta: el contraste adicional entre 3 y 4+ visitas es pequeño descriptivamente, por lo que la evidencia respalda con mayor claridad una asociación entre utilización del servicio y desempeño que un beneficio incremental específico de exceder el umbral PPA.

El análisis temporal muestra que el comportamiento no parece dominado por estudiantes que realizan tres visitas en un solo día: solo 1.48\% de quienes alcanzan PPA concentra sus primeras tres visitas en la misma fecha. A nivel agregado, la demanda presenta una periodicidad aproximadamente mensual: la mediana entre picos detectados es 30.5 días y la máxima autocorrelación agrupada aparece a 30 días. Este patrón es compatible con el ciclo institucional aproximado de cuatro evaluaciones por semestre, aunque las fechas de examen por profesor no están observadas y no puede atribuirse causalmente cada pico a una evaluación. La regularidad temporal individual, una vez controlado el número de visitas, no muestra una asociación adicional clara con desempeño ($\hat\beta=-0.136$, $p=0.413$). Finalmente, entre 4,211 estudiantes que progresan posteriormente a Cálculo I con cobertura de asesorías, la probabilidad de volver a usar el CMAT aumenta de 16.9\% entre no usuarios iniciales a 60.6\% entre quienes habían realizado 4 o más visitas en Matemáticas Universitarias.

En conjunto, los resultados describen una asociación robusta entre uso del CMAT y desempeño, una fuerte persistencia conductual posterior y un patrón temporal compatible con el ritmo de evaluación. No establecen causalidad: la ausencia de medidas individuales pretratamiento impide descartar autoselección por motivación, preparación previa, organización o necesidad académica.
\end{abstract}

\noindent\textbf{Palabras clave:} apoyo matemático universitario; help-seeking; tutoría; desempeño académico; PPA; autoselección; persistencia; análisis longitudinal.

\tableofcontents
\listoffigures
\listoftables
\clearpage
\pagenumbering{arabic}

\section{Introducción}

Los centros universitarios de apoyo matemático buscan complementar la enseñanza formal mediante asesorías voluntarias, resolución de dudas y acompañamiento académico. La literatura internacional sobre \emph{mathematics and statistics support} ha crecido considerablemente y, en términos generales, documenta asociaciones favorables entre utilización de estos servicios y desempeño; al mismo tiempo, las revisiones recientes subrayan que gran parte de la evidencia está limitada por autoselección, heterogeneidad en la definición de asistencia y diseños observacionales difíciles de comparar \citep{mullen2024,lawson2020}. Estudios directamente relacionados con calificaciones de cursos matemáticos han encontrado asociaciones positivas entre asistencia y desempeño \citep{rickard2018,macanbhaird2009,jacob2019}, mientras que diseños más rigurosos han mostrado resultados más matizados y recuerdan que una asociación favorable no equivale a un efecto causal \citep{buchele2024}.

El contexto del CMAT permite estudiar una dimensión adicional. Durante el primer intento de Matemáticas Universitarias, los estudiantes pueden obtener un punto de participación institucional (PPA) al alcanzar tres visitas cualesquiera al CMAT durante el semestre. Este punto no modifica la calificación de la materia; funciona como incentivo externo para participar en actividades universitarias de primer año. Por tanto, el umbral $V=3$ tiene significado administrativo, mientras que $V>3$ representa utilización observada después de haber satisfecho el requisito. Incentivos de baja intensidad y diseños de \emph{encouragement} han sido utilizados en educación superior para estudiar adopción de servicios de apoyo y hábitos de estudio \citep{paloyo2016,pugatch2018}.

El propósito de este documento es presentar de forma integrada los resultados del nuevo diseño analítico. No se limita a reportar pruebas estadísticas: para cada pregunta se explicita qué se pregunta, qué responde la evidencia, qué puede deducirse razonablemente, qué no puede deducirse y qué mecanismos alternativos podrían producir el patrón observado.

\begin{cautionbox}
El estudio es observacional. Las expresiones ``asociación'', ``diferencia'', ``persistencia'' y ``patrón compatible con'' se utilizan deliberadamente. No se interpreta que una visita al CMAT cause automáticamente un aumento de calificación. Las decisiones de acudir al servicio pueden depender de variables no observadas, como motivación, preparación matemática previa, disciplina, percepción de dificultad o disponibilidad de tiempo.
\end{cautionbox}

\section{Preguntas de investigación y lógica del estudio}

La contribución del proyecto se organiza alrededor de cinco preguntas relacionadas pero conceptualmente distintas.

\begin{table}[H]
\centering
\caption{Mapa de preguntas de investigación.}
\label{tab:rqmap}
\begin{tabularx}{\textwidth}{@{}p{1.25cm}X X@{}}
\toprule
\textbf{RQ} & \textbf{Pregunta} & \textbf{Resultado que la responde} \\
\midrule
RQ1a & ¿Cómo se distribuye el uso alrededor del umbral PPA de tres visitas? & Distribución exacta de visitas, proporción en $V=3$, continuación $P(V\geq k+1\mid V\geq k)$ y comparación descriptiva con Cálculo I. \\
RQ1b & ¿Cuándo se utiliza el CMAT y existe concentración temporal compatible con el ciclo de evaluaciones? & Visitas por día, picos por periodo, separación entre picos, autocorrelación y concentración de tres visitas en un mismo día. \\
RQ2a & ¿El uso del CMAT se asocia con mayor desempeño continuo en el primer intento de MU? & Medias por grupo, pruebas robustas, efectos fijos profesor-periodo, dosis por grupos y sensibilidades. \\
RQ2b & A igual intensidad aproximada, ¿una utilización más distribuida en el tiempo se asocia con mayor desempeño? & Métrica de regularidad mensual y modelos que controlan el número total de visitas. \\
RQ3 & ¿El patrón de uso persiste posteriormente en Cálculo I, donde el PPA ya no aplica? & Probabilidad de cualquier visita en Cálculo y número posterior de visitas, condicionado a progresar al curso. \\
\bottomrule
\end{tabularx}
\end{table}

La secuencia conceptual puede expresarse como
\[
\begin{aligned}
\text{PPA}&\longrightarrow\text{adopción del CMAT}\longrightarrow\text{uso más allá del umbral}\\
&\longleftrightarrow\text{desempeño}\longrightarrow\text{trayectoria posterior}.
\end{aligned}
\]
pero las flechas no deben interpretarse causalmente. Desde la perspectiva de \emph{academic help-seeking}, acudir a asesorías puede ser una conducta autorregulatoria, pero la misma conducta está relacionada con características del estudiante que también predicen desempeño \citep{fong2023}.

\section{Datos y diseño analítico}

\subsection{Cohorte primaria y reconstrucción temporal}

La unidad primaria es el estudiante en su \textbf{primer intento estándar} de Matemáticas Universitarias. Los intentos posteriores se excluyen, pero un estudiante que posteriormente repite conserva su primer intento. Los códigos EQV, REV y AC se eliminan porque corresponden a acreditaciones externas o equivalentes cuyo profesor y aula no están observados. Los estados BV, RT y BA se mantienen como resultados académicos adversos.

La construcción de la cohorte se resume en la Tabla~\ref{tab:cohort}.

\begin{table}[H]
\centering
\caption{Flujo de construcción de las cohortes.}
\label{tab:cohort}
\begin{tabular}{@{}lr@{}}
\toprule
\textbf{Etapa} & \textbf{$N$} \\
\midrule
Primeros intentos elegibles de MU, todos los periodos & 7,777 \\
Cohorte primaria de MU con cobertura de asesorías & 6,627 \\
Primeros intentos de Cálculo I con cobertura, comparador descriptivo & 5,046 \\
Estudiantes que progresan de MU a Cálculo posterior & 5,048 \\
Pares longitudinales con cobertura de asesorías en Cálculo & 4,211 \\
\bottomrule
\end{tabular}
\end{table}

La exposición primaria para el estudiante $i$ es
\[
\begin{aligned}
V_i=\#\{\,&\text{visitas cualesquiera al CMAT durante el mismo año y periodo}\\
&\text{del primer intento de Matemáticas Universitarias}\}.
\end{aligned}
\]
Esta definición corrige un problema del análisis exploratorio previo, en el cual las visitas podían acumularse a lo largo de varios años y asignarse posteriormente a observaciones académicas de periodos diferentes.

\subsection{Grupos de exposición y significado del PPA}

Se construyen cuatro grupos:
\[
G_i=\begin{cases}
0, & V_i=0,\\
1, & V_i\in\{1,2\},\\
2, & V_i=3,\\
3, & V_i\geq4.
\end{cases}
\]
Además,
\[
\text{PPA\_REACHED}_i=1(V_i\geq3),\qquad
\text{BEYOND\_PPA}_i=1(V_i>3).
\]
El grupo $V>3$ no se interpreta como ``motivación intrínseca''; únicamente describe uso que continúa después del mínimo requerido.

\subsection{Outcome primario: desempeño continuo}

La respuesta principal es la calificación final estandarizada dentro de un contexto profesor-periodo:
\[
Z_i=\frac{Y_i-\mu_{p,t}}{\sigma_{p,t}},
\]
donde $Y_i$ es la calificación del estudiante, $p$ el profesor y $t$ el periodo. Esta transformación reduce la confusión producida por diferencias persistentes de nivel y dispersión entre contextos de evaluación.

Para BV, RT y BA se asume institucionalmente que el desempeño subyacente es inferior a 7.5. El análisis primario continuo utiliza imputación adversa por debajo de ese umbral y después estandariza. Se preespecifican dos sensibilidades: imputación uniforme bajo 7.5 y análisis de casos con calificación numérica observada.

El outcome secundario es aprobación:
\[
PASS_i=1(Y_i\geq7.5),
\]
con BV, RT y BA clasificados como no aprobación. Este outcome es importante porque no requiere asignar una calificación numérica a las bajas.

\subsection{Modelos y alcance inferencial}

La especificación principal para el contraste $V>3$ es
\[
Z_i=\beta\,1(V_i>3)+\alpha_{p,t}+\gamma_{c(i)}+\varepsilon_i,
\]
donde $\alpha_{p,t}$ son efectos fijos profesor-periodo y $\gamma_{c(i)}$ representa carrera. Los errores estándar se agrupan por profesor-periodo. Además se reportan pruebas de Welch, Mann--Whitney, Brunner--Munzel, probabilidad de superioridad y delta de Cliff.

Se calcula también una sensibilidad con ponderación por propensity score basada en carrera, año, periodo y profesor. Los métodos de propensity score pueden mejorar balance sobre variables observadas \citep{rosenbaum1983,austin2011}, pero no eliminan confusión por variables individuales no observadas \citep{hernan2020}. En este dataset no existen diagnóstico matemático, promedio de preparatoria ni otras covariables individuales pretratamiento.

\section{Resultados}

\subsection{RQ1a: uso alrededor del umbral PPA}

\begin{questionbox}
¿La obligación de alcanzar tres visitas parece producir una acumulación mecánica de estudiantes exactamente en $V=3$? ¿Y qué ocurre con la probabilidad de continuar después de alcanzar el umbral?
\end{questionbox}

De 6,627 estudiantes, 5,393 (81.4\%) no registraron ninguna visita durante el semestre de MU. Hubo 517 con una visita, 245 con dos, 170 con exactamente tres y 302 con cuatro o más. En total, 1,234 estudiantes (18.6\%) utilizaron el CMAT al menos una vez; 472 (7.1\% de la cohorte y 38.2\% de los usuarios) alcanzaron el PPA y 302 (4.6\% de la cohorte) continuaron más allá del umbral.

\begin{figure}[H]
\centering
\includegraphics[width=0.92\textwidth]{figures/01_visit_distribution_mu_vs_calculus.png}
\caption{Distribución de visitas al CMAT durante el periodo académico para Matemáticas Universitarias y Cálculo I. La línea vertical indica el umbral PPA en tres visitas.}
\label{fig:visitdist}
\end{figure}

La Figura~\ref{fig:visitdist} muestra dos hechos. Primero, la distribución está fuertemente concentrada en cero y en conteos bajos. Segundo, no aparece una acumulación excepcional en $V=3$. El índice local de \emph{bunching}, definido como
\[
B=\frac{P(V=3)}{[P(V=2)+P(V=4)]/2},
\]
es $B=0.997$ en MU. Un valor cercano a uno significa que la masa en tres visitas es aproximadamente la esperada por interpolación local entre dos y cuatro, no un pico extraordinario.

La comparación con Cálculo I refuerza esta lectura descriptiva: aunque en Cálculo no existe PPA, el 22.0\% registra al menos una visita y la media de visitas es 0.667, frente a 18.6\% y 0.539 en MU. El índice local de bunching de Cálculo también es cercano a uno ($0.951$). Esta comparación no identifica el efecto del PPA porque cambian dificultad, composición y semestre; sí muestra que el uso del CMAT no desaparece cuando el incentivo deja de estar presente.

\begin{figure}[H]
\centering
\includegraphics[width=0.92\textwidth]{figures/02_continuation_probability.png}
\caption{Probabilidad de continuar a una visita adicional, $P(V\geq k+1\mid V\geq k)$, para MU y Cálculo I.}
\label{fig:continuation}
\end{figure}

Entre quienes alcanzaron tres visitas, 63.98\% continuaron a una cuarta. En términos simples, casi dos tercios de quienes obtienen el PPA no se detienen exactamente al cumplir el requisito. La probabilidad de continuación no presenta una caída abrupta en $k=3$: en MU pasa de 65.8\% al ir de dos a tres visitas a 64.0\% al ir de tres a cuatro.

\begin{answerbox}
La evidencia no muestra un pico especial en exactamente tres visitas ni un abandono masivo inmediatamente después del PPA. De los 472 estudiantes que alcanzan tres visitas, 302 continúan a cuatro o más.
\end{answerbox}

\begin{interpretbox}
El PPA probablemente funciona como una puerta de entrada para algunos estudiantes, pero el patrón observado no es consistente con una historia puramente mecánica de ``voy exactamente tres veces y dejo de asistir''. También es posible que los estudiantes que alcanzan tres sean ya una población más proclive a buscar ayuda. La comparación con Cálculo, donde el uso es incluso más frecuente pese a no existir PPA, sugiere que dificultad académica, familiaridad con el servicio y una propensión estable a pedir ayuda pueden ser mecanismos importantes.
\end{interpretbox}

\subsection{RQ1b: cuándo se utiliza el CMAT}

\begin{questionbox}
¿Los estudiantes concentran las tres visitas del PPA en una sola fecha? ¿Y la demanda agregada del CMAT presenta un ciclo temporal compatible con aproximadamente cuatro evaluaciones mensuales por semestre?
\end{questionbox}

\subsubsection{Concentración de visitas en un mismo día}

Entre los 1,234 usuarios de la cohorte de MU, 155 tuvieron al menos un día con dos o más visitas y 24 tuvieron algún día con tres o más. Restringiendo a los 472 que alcanzaron PPA, 133 (28.2\%) registraron algún día con dos o más visitas, 24 (5.1\%) algún día con tres o más, y solo 7 (1.48\%) realizaron sus \emph{primeras tres visitas} en la misma fecha.

\begin{figure}[H]
\centering
\includegraphics[width=0.82\textwidth]{figures/07_same_day_ppa_behavior.png}
\caption{Máximo número de visitas en un mismo día entre quienes alcanzaron el PPA. El 71.8\% nunca tuvo más de una visita diaria.}
\label{fig:sameday}
\end{figure}

\begin{table}[H]
\centering
\caption{Concentración intradía entre estudiantes que alcanzaron el PPA ($N=472$).}
\label{tab:sameday}
\begin{tabular}{@{}lrr@{}}
\toprule
\textbf{Indicador} & \textbf{Estudiantes} & \textbf{Proporción} \\
\midrule
Algún día con 2+ visitas & 133 & 28.2\% \\
Algún día con 3+ visitas & 24 & 5.1\% \\
Primeras tres visitas el mismo día & 7 & 1.48\% \\
Día de completar PPA con 2+ visitas & 68 & 14.4\% \\
Día de completar PPA con 3+ visitas & 18 & 3.8\% \\
\bottomrule
\end{tabular}
\end{table}

\begin{answerbox}
La estrategia extrema de acumular las tres visitas necesarias en un solo día existe, pero es rara. La gran mayoría de quienes alcanzan el PPA distribuye las visitas en más de una fecha.
\end{answerbox}

\begin{interpretbox}
Este resultado ayuda a descartar una explicación trivial del uso observado: el grupo $V\geq3$ no está compuesto principalmente por estudiantes que ``fabrican'' tres registros en una sola jornada. No obstante, no sabemos la motivación individual; distribuir tres visitas en días distintos tampoco demuestra interés académico intrínseco.
\end{interpretbox}

\subsubsection{Ciclo de demanda aproximadamente mensual}

Las fechas exactas de los exámenes dependen del profesor y no están disponibles. Institucionalmente, la mayoría de los cursos realiza aproximadamente cuatro evaluaciones, una por mes, con fechas de inicio y fin de semestre comunes. Por ello, la inferencia posible es sobre \emph{periodicidad}, no sobre coincidencia con un examen específico.

\begin{figure}[H]
\centering
\includegraphics[width=0.98\textwidth]{figures/06_term_temporal_peaks.png}
\caption{Actividad diaria del CMAT por periodo y picos detectados sobre una ventana móvil de siete días. La repetición aproximada de máximos a intervalos mensuales es visible en varios periodos.}
\label{fig:termpeaks}
\end{figure}

En 11 periodos analizados, el algoritmo detecta una mediana de 5 picos por periodo en toda la población del CMAT. Entre 38 intervalos consecutivos entre picos, la mediana es 30.5 días, con rango intercuartil 28.0--33.75 días. Veintiocho de los 38 intervalos (73.7\%) caen entre 24 y 38 días.

\begin{figure}[H]
\centering
\begin{subfigure}[t]{0.49\textwidth}
\centering
\includegraphics[width=\linewidth]{figures/08_peak_spacing_days.png}
\caption{Separación entre picos detectados.}
\end{subfigure}
\hfill
\begin{subfigure}[t]{0.49\textwidth}
\centering
\includegraphics[width=\linewidth]{figures/09_monthly_periodicity_acf.png}
\caption{Autocorrelación de la serie diaria ajustada.}
\end{subfigure}
\caption{Dos diagnósticos independientes de periodicidad temporal. La máxima autocorrelación agrupada se observa a 30 días ($\rho\approx0.210$).}
\label{fig:periodicity}
\end{figure}

Después de remover el promedio por día de la semana y una tendencia lineal, la autocorrelación agrupada alcanza su máximo a 30 días, $\rho(30)=0.210$. El periodograma restringido a ciclos de 21--42 días produce una mediana de periodo dominante de aproximadamente 27 días entre los periodos con longitud suficiente. Aunque la señal no es idéntica en todos los semestres, tres diagnósticos diferentes --picos, autocorrelación y frecuencia-- apuntan a una escala temporal cercana a un mes.

\begin{answerbox}
La demanda presenta una estructura temporal recurrente cercana a cuatro semanas. La mediana entre picos es 30.5 días y la máxima autocorrelación agrupada aparece a 30 días.
\end{answerbox}

\begin{interpretbox}
El patrón es coherente con un calendario en el que la mayoría de los profesores evalúa aproximadamente una vez al mes: estudiantes de cursos distintos pueden tener parciales en días diferentes, pero al agregarlos aparece una ``onda'' de demanda alrededor de cada ciclo. Esta explicación es plausible, no demostrada. También podrían contribuir fechas de entrega, semanas de mayor carga académica, periodos de bajas o campañas internas del CMAT. La interpretación correcta es ``compatible con el ritmo de evaluación'', no ``causado por los exámenes''.
\end{interpretbox}

\subsection{RQ2a: utilización y desempeño continuo}

\begin{questionbox}
Entre estudiantes que cursan MU por primera vez, ¿quienes utilizan más el CMAT se ubican en mejores posiciones relativas dentro de su contexto profesor-periodo?
\end{questionbox}

La Figura~\ref{fig:primaryoutcome} presenta el resultado descriptivo central.

\begin{figure}[H]
\centering
\includegraphics[width=0.88\textwidth]{figures/03_primary_outcome_by_visit_group.png}
\caption{Desempeño estandarizado medio por grupo de visitas, con intervalos de confianza al 95\%.}
\label{fig:primaryoutcome}
\end{figure}

\begin{table}[H]
\centering
\caption{Outcome continuo por grupo de visitas.}
\label{tab:primarydesc}
\begin{tabular}{@{}lrrrr@{}}
\toprule
\textbf{Visitas} & \textbf{$N$} & \textbf{Media $Z$} & \textbf{Mediana $Z$} & \textbf{\CI\ de la media} \\
\midrule
0 & 5,393 & -0.060 & 0.315 & $[-0.087,-0.033]$ \\
1--2 & 762 & 0.218 & 0.453 & $[0.160,0.276]$ \\
3 & 170 & 0.325 & 0.543 & $[0.216,0.434]$ \\
4+ & 302 & 0.342 & 0.535 & $[0.262,0.423]$ \\
\bottomrule
\end{tabular}
\end{table}

El patrón es monotónico en la media, pero no lineal. El salto principal ocurre al pasar de cero visitas a alguna utilización. Los coeficientes del modelo de dosis, siempre respecto del grupo de cero visitas y con efectos fijos profesor-periodo, son:
\[
\hat\beta_{1-2}=0.318,\qquad
\hat\beta_{3}=0.443,\qquad
\hat\beta_{4+}=0.468.
\]
La diferencia entre los coeficientes de 3 y 4+ es pequeña en comparación con el salto inicial. Esto sugiere una posible meseta descriptiva después de tres visitas; el análisis disponible no prueba un retorno marginal positivo de cada visita adicional.

\subsubsection{Contraste principal $V>3$ frente a $V\leq3$}

El grupo con cuatro o más visitas tiene media $0.342$ frente a $-0.016$ en el grupo con tres o menos. La diferencia es
\[
\Delta\bar Z=0.359,\qquad \CI=[0.277,0.439],
\]
con Welch $t=8.37$ y $p=1.32\times10^{-15}$. La diferencia de medianas es $0.198$ (\CI\ $[0.115,0.241]$). Mann--Whitney y Brunner--Munzel también rechazan igualdad distributiva ($p<10^{-11}$).

La probabilidad de superioridad es
\[
CL=0.616,
\]
lo que significa que, al seleccionar aleatoriamente una observación de cada grupo, la del grupo $V>3$ tiene una probabilidad aproximada de 61.6\% de tener mayor $Z$. El delta de Cliff es $0.233$, una separación positiva de magnitud no trivial.

\begin{table}[H]
\centering
\caption{Modelos principales con efectos fijos profesor-periodo.}
\label{tab:femodels}
\begin{tabular}{@{}lrrrr@{}}
\toprule
\textbf{Modelo} & \textbf{Estimación} & \textbf{EE} & \textbf{\CI} & \textbf{$p$} \\
\midrule
Profesor-periodo FE & 0.388 & 0.060 & $[0.270,0.506]$ & $1.10\times10^{-10}$ \\
Profesor-periodo FE + carrera & 0.375 & 0.060 & $[0.257,0.493]$ & $4.08\times10^{-10}$ \\
\bottomrule
\end{tabular}
\end{table}

El ajuste por contexto de evaluación y carrera reduce poco la estimación. Una sensibilidad de propensity score basada únicamente en carrera, año, periodo y profesor produce un efecto ATT ponderado de $0.373$ y reduce el máximo desequilibrio estandarizado observado de 0.443 a 0.116. Esta estabilidad es útil como prueba de robustez sobre variables disponibles, pero no resuelve autoselección individual.

\begin{answerbox}
La utilización del CMAT está fuertemente asociada con mejor desempeño relativo. El resultado aparece en medias, medianas, pruebas robustas, modelos profesor-periodo y ajuste por carrera. Sin embargo, la desagregación en 0, 1--2, 3 y 4+ muestra que la evidencia más clara es ``usar el servicio versus no usarlo o usarlo poco'', no un beneficio incremental grande de 4+ frente a exactamente 3.
\end{answerbox}

\begin{interpretbox}
Hay al menos cuatro explicaciones compatibles con los datos. (1) Las asesorías pueden mejorar comprensión y desempeño. (2) Estudiantes más organizados o motivados pueden ser simultáneamente más proclives a asistir y a obtener mejores notas. (3) Estudiantes con mayor dificultad pueden asistir más, generando confusión en sentido contrario y atenuando una posible asociación favorable. (4) Una o pocas visitas podrían resolver dudas importantes, de modo que los retornos adicionales sean decrecientes. La meseta entre 3 y 4+ es consistente con (4), pero también con mezcla de perfiles de necesidad; no permite elegir entre mecanismos.
\end{interpretbox}

\subsubsection{Sensibilidad al tratamiento de BV, RT y BA}

La imputación uniforme bajo 7.5 produce una diferencia media de $0.343$ y Cliff $\delta=0.227$, muy cercana al análisis primario. En cambio, al analizar únicamente estudiantes con calificación numérica, la diferencia media se reduce a $0.211$ y Cliff $\delta=0.075$.

Esta atenuación es substantivamente importante. Parte de la separación entre grupos está asociada a la presencia de resultados adversos no numéricos. Esto no significa que la imputación sea un artefacto necesariamente: evitar una baja o un retiro podría ser un resultado académico relevante. Pero sí implica que la magnitud del outcome continuo depende parcialmente de cómo se incorpora ese tipo de resultado y que el outcome secundario de aprobación es necesario como verificación independiente.

\subsubsection{Outcome secundario: aprobación}

Las tasas de aprobación son 72.1\% para cero visitas, 81.0\% para 1--2, 81.8\% para exactamente 3 y 86.4\% para 4+. El modelo lineal de probabilidad con efectos fijos profesor-periodo y carrera estima para $V>3$ una diferencia de 0.192 puntos de probabilidad (\CI\ $[0.138,0.247]$, $p=5.59\times10^{-12}$).

\begin{cautionbox}
La magnitud del modelo de aprobación es secundaria y no debe interpretarse como ``el CMAT aumenta 19 puntos porcentuales la probabilidad de aprobar''. Es una diferencia ajustada observacional. Variables individuales no observadas podrían explicar parte sustancial de ella.
\end{cautionbox}

\subsection{RQ2b: regularidad temporal y desempeño}

\begin{questionbox}
A igual número aproximado de visitas, ¿es mejor distribuirlas a lo largo del semestre que concentrarlas en pocos meses?
\end{questionbox}

Se define
\[
R_i=\frac{\min(M_i,4)}{\min(V_i,4)},
\]
donde $M_i$ es el número de meses calendario distintos con al menos una visita. Entre estudiantes que alcanzan PPA, la cantidad de visitas y la dispersión temporal están fuertemente relacionadas: quienes aparecen en un solo mes tienen 3.53 visitas en promedio; quienes aparecen en cuatro meses, 9.07. Por ello no basta con comparar $R_i$ sin controlar intensidad.

La especificación primaria se restringe a estudiantes con $V\geq3$ y calificación numérica final ($N=413$), controla conteos de visitas 3, 4, 5, 6, 7 y 8+, efectos fijos profesor-periodo y carrera. El coeficiente de regularidad es
\[
\hat\beta_R=-0.136,\qquad \CI=[-0.462,0.190],\qquad p=0.413.
\]
Una sensibilidad con outcomes adversos imputados da $-0.166$ ($p=0.489$), y el número efectivo de semanas por visita da $-0.049$ ($p=0.800$).

\begin{figure}[H]
\centering
\includegraphics[width=0.80\textwidth]{figures/10_exact3_regularity_vs_performance.png}
\caption{Análisis exploratorio entre estudiantes con exactamente tres visitas: desempeño según número de meses calendario activos.}
\label{fig:regularity}
\end{figure}

Entre quienes tienen exactamente tres visitas, la media de $Z$ es 0.291 si las visitas se concentran en un mes, 0.291 en dos meses y 0.403 en tres meses. La diferencia visual del tercer grupo es sugestiva pero los intervalos se superponen y el tamaño muestral es reducido (26, 73 y 45 estudiantes, respectivamente).

\begin{answerbox}
No se detecta evidencia clara de que la regularidad temporal aporte información adicional sobre desempeño una vez controlada la intensidad de uso. La hipótesis preespecificada $\beta_R>0$ no recibe apoyo estadístico con la métrica actual.
\end{answerbox}

\begin{interpretbox}
Un resultado nulo también es informativo. Puede significar que la cantidad de exposición es más relevante que su dispersión temporal, pero también puede reflejar potencia limitada, una métrica mensual demasiado gruesa o heterogeneidad en las fechas de evaluación. Además, concentrar visitas antes de un examen puede ser una estrategia racional, no necesariamente una conducta académica inferior. Por tanto, el resultado no autoriza afirmar que ``la regularidad no importa''; solo indica que no se observa una asociación independiente clara con los datos y definición actuales.
\end{interpretbox}

\subsection{RQ3: persistencia hacia Cálculo I}

\begin{questionbox}
Entre estudiantes que posteriormente progresan a Cálculo I, ¿el uso previo del CMAT durante MU predice que vuelvan a utilizar el servicio cuando el incentivo PPA ya no está disponible?
\end{questionbox}

La cohorte longitudinal contiene 4,211 estudiantes con cobertura de asesorías durante su primer Cálculo posterior. De ellos, 3,842 cursan Cálculo en el semestre académico inmediatamente siguiente; el resto presenta distintos rezagos.

\begin{figure}[H]
\centering
\includegraphics[width=0.86\textwidth]{figures/04_longitudinal_persistence_to_calculus.png}
\caption{Probabilidad de registrar al menos una visita al CMAT durante Cálculo I según patrón de uso previo en Matemáticas Universitarias.}
\label{fig:persistence}
\end{figure}

\begin{table}[H]
\centering
\caption{Persistencia de utilización en Cálculo I.}
\label{tab:persistence}
\begin{tabular}{@{}lrrrr@{}}
\toprule
\textbf{Grupo MU} & \textbf{$N$} & \textbf{Alguna visita en Cálculo} & \textbf{Media visitas Cálculo} & \textbf{Mediana} \\
\midrule
0 & 3,274 & 16.9\% & 0.43 & 0 \\
1--2 & 567 & 41.1\% & 1.28 & 0 \\
3 & 129 & 51.2\% & 1.78 & 1 \\
4+ & 241 & 60.6\% & 3.12 & 1 \\
\bottomrule
\end{tabular}
\end{table}

La relación es monotónica: cuanto mayor es la utilización durante MU, mayor es la probabilidad y la intensidad de utilización posterior. Un estudiante del grupo 4+ tiene una probabilidad observada de volver al CMAT en Cálculo aproximadamente 3.6 veces la del grupo que no había acudido en MU ($60.6/16.9\approx3.59$).

\begin{answerbox}
El uso del CMAT muestra una fuerte persistencia longitudinal. La probabilidad de volver a utilizar el servicio en Cálculo aumenta sistemáticamente con el nivel de uso previo en MU, incluso cuando el PPA ya no aplica.
\end{answerbox}

\begin{interpretbox}
Existen al menos dos mecanismos principales. Uno es aprendizaje conductual: el PPA o el primer contacto puede reducir barreras, familiarizar al estudiante con el centro y convertir la búsqueda de ayuda en una estrategia disponible. El otro es selección estable: estudiantes con mayor propensión a pedir ayuda, mayor disciplina o necesidades académicas recurrentes podrían utilizar el centro en ambos semestres sin que la primera experiencia cause la segunda. Con los datos actuales no es posible separar estos mecanismos. No obstante, la persistencia es un hallazgo sustantivo porque muestra que el comportamiento no termina con el incentivo inicial.
\end{interpretbox}

\section{Síntesis de respuestas y deducciones}

\begin{table}[H]
\centering
\caption{Resumen integrado de hallazgos.}
\label{tab:summary}
\begin{tabularx}{\textwidth}{@{}p{1.25cm}X X@{}}
\toprule
\textbf{RQ} & \textbf{Respuesta empírica} & \textbf{Deducción prudente} \\
\midrule
RQ1a & No hay bunching claro en $V=3$; 64\% de quienes llegan a 3 continúan a 4+. & El PPA no parece generar un patrón dominante de cumplimiento exacto y abandono inmediato. \\
RQ1b & Primeras tres visitas el mismo día: 1.48\%; mediana entre picos: 30.5 días; ACF máxima a 30 días. & El uso está distribuido en el tiempo y la demanda agregada tiene un ciclo compatible con evaluaciones aproximadamente mensuales. \\
RQ2a & $V>3$ se asocia con $+0.375Z$ ajustado; gradiente 0 $<$ 1--2 $<$ 3 $\approx$ 4+. & Hay asociación robusta entre uso y desempeño; la evidencia de beneficio incremental después de 3 es mucho más débil que la diferencia entre no uso y uso. \\
RQ2b & $\hat\beta_R=-0.136$, $p=0.413$. & No se detecta un valor predictivo adicional claro de distribuir visitas entre más meses una vez controlada intensidad. \\
RQ3 & Visita posterior en Cálculo: 16.9\%, 41.1\%, 51.2\%, 60.6\% según grupo previo. & Existe una fuerte persistencia de help-seeking, compatible tanto con aprendizaje conductual como con heterogeneidad estable entre estudiantes. \\
\bottomrule
\end{tabularx}
\end{table}

\subsection{Qué parece ser la historia principal de los datos}

Los resultados son más coherentes con una historia de \textbf{adopción y persistencia} que con una historia de simple cumplimiento administrativo. Tres observaciones sostienen esta lectura. Primero, no aparece una acumulación extraordinaria exactamente en tres visitas. Segundo, casi dos tercios de quienes alcanzan el umbral continúan a una cuarta visita. Tercero, la utilización previa predice fuertemente la utilización posterior en Cálculo, donde el PPA ya no aplica.

Sin embargo, el outcome de desempeño obliga a una interpretación más cuidadosa. El contraste binario $V>3$ frente a $V\leq3$ es fuerte, pero el modelo por grupos muestra que $V=3$ y $V=4+$ tienen niveles de desempeño muy parecidos. Por tanto, el mensaje científico más sólido no es ``más allá de tres visitas produce una ganancia adicional grande'', sino:
\begin{tcolorbox}[colback=cmatlight,colframe=cmatblue,boxrule=0.6pt]
\centering
\textbf{La utilización del CMAT está asociada con mejor desempeño relativo,}\par
\textbf{con una separación observada que parece aplanarse alrededor de tres visitas.}
\end{tcolorbox}
La palabra ``rendimientos'' aquí es descriptiva, no causal.

\subsection{Tesis, explicación rival y síntesis}

\paragraph{Tesis.} Las asesorías proporcionan apoyo académico útil, reducen dudas y ayudan a los estudiantes a ubicarse en posiciones más altas dentro de su contexto de evaluación. La consistencia entre outcome continuo, aprobación, pruebas robustas y seguimiento longitudinal es compatible con esta interpretación.

\paragraph{Explicación rival.} Los usuarios del CMAT son distintos antes de acudir. Podrían ser más disciplinados, más conscientes de sus dificultades o más motivados. También puede ocurrir selección por necesidad: estudiantes con mayor dificultad buscan más asesorías, lo que puede ocultar parte de un posible beneficio. Sin covariables individuales pretratamiento, ambos tipos de selección permanecen abiertos.

\paragraph{Síntesis.} La evidencia actual justifica afirmar que el CMAT identifica y acompaña una población cuyo comportamiento académico y de búsqueda de ayuda es sistemáticamente distinto. La estabilidad de la asociación tras controlar profesor-periodo, carrera y variables administrativas observadas hace menos probable que el resultado se deba únicamente a heterogeneidad de aula; no elimina heterogeneidad individual. En términos de publicación, este es un estudio observacional longitudinal con mecanismos plausibles y limitaciones explícitas, no una evaluación causal definitiva.

\section{Implicaciones institucionales}

Los resultados sugieren tres aplicaciones inmediatas para el CMAT.

\begin{enumerate}[leftmargin=*]
\item \textbf{Planeación de capacidad.} La periodicidad cercana a 30 días permite anticipar ventanas de mayor demanda. Aunque las fechas de examen varíen entre profesores, la agregación institucional genera ciclos suficientemente visibles para planear disponibilidad de asesores.
\item \textbf{El PPA como mecanismo de exposición.} El hecho de que muchos estudiantes continúen después de la tercera visita y vuelvan en Cálculo sugiere que el primer contacto puede tener valor más allá del cumplimiento administrativo. Esto debe estudiarse como mecanismo de adopción, no asumirse causalmente.
\item \textbf{Seguimiento longitudinal.} Los estudiantes con uso recurrente constituyen un grupo identificable para estudiar trayectorias posteriores, necesidades de apoyo y potenciales intervenciones de seguimiento.
\end{enumerate}

\section{Limitaciones}

\begin{enumerate}[leftmargin=*]
\item \textbf{Autoselección no observada.} No existen medidas individuales previas a MU como diagnóstico matemático, promedio de preparatoria o examen de admisión. Esta es la limitación inferencial principal.
\item \textbf{Bajas sin fecha exacta.} BV, RT y BA se conocen por periodo, pero no por fecha. Esto limita el análisis de exposición temporal y explica por qué RQ2b usa casos con calificación numérica como población primaria.
\item \textbf{Imputación del outcome continuo.} El análisis completo muestra que la magnitud de la asociación se atenúa cuando se excluyen resultados adversos no numéricos. El outcome de aprobación ayuda a mostrar que el patrón no depende únicamente de asignarles un número.
\item \textbf{Fechas de examen no observadas.} El ciclo de 30 días es compatible con evaluaciones aproximadamente mensuales, pero no puede vincularse a parciales específicos por profesor.
\item \textbf{Comparación MU--Cálculo no causal.} Los estudiantes que llegan a Cálculo son una población seleccionada por progresión académica y el curso tiene distinta dificultad; no es un grupo control del PPA.
\item \textbf{Cobertura de asesorías.} La inferencia primaria se restringe a periodos con registros efectivos de asesoría; periodos académicos sin cobertura no se interpretan como cero visitas.
\end{enumerate}

\section{Próximos análisis recomendados}

La siguiente fase debería priorizar profundidad antes que añadir pruebas indiscriminadamente.

\begin{enumerate}[leftmargin=*]
\item Estimar modelos de dosis no lineales o splines para evaluar formalmente si existe meseta después de 2--3 visitas.
\item Comparar directamente $V=3$ frente a $V=4+$ dentro de contextos profesor-periodo, evitando que el gran grupo $V=0$ domine el contraste binario principal.
\item Modelar la probabilidad de persistencia en Cálculo con ajuste por desempeño de MU y semestre de progresión, dejando claro que condicionar por desempeño cambia el estimando.
\item Explorar heterogeneidad del efecto asociativo por profesor, carrera y periodo únicamente con hipótesis preespecificadas y corrección por múltiples comparaciones.
\item Si en el futuro se recuperan covariables pretratamiento, reestimar con métodos doblemente robustos o ponderación de propensity score más informativa.
\item Para RQ1b, construir modelos de conteo diarios con efectos de día de semana y componentes armónicos de 28--31 días para cuantificar la periodicidad en lugar de depender solo de detección de picos.
\end{enumerate}

\section{Conclusión}

El rediseño del estudio modifica de manera importante la interpretación del reporte exploratorio original. Al restringir el análisis al primer intento de una misma materia, enlazar las visitas exclusivamente con el semestre correspondiente y estandarizar el desempeño dentro de profesor-periodo, se obtiene una comparación mucho más coherente temporal y académicamente.

La asociación principal permanece: los estudiantes que utilizan más el CMAT se ubican, en promedio, en posiciones relativas superiores y presentan mayores tasas de aprobación. La evidencia no muestra un pico artificial exactamente en tres visitas; la mayoría de quienes alcanzan el PPA distribuye sus visitas entre distintas fechas y 64\% continúa más allá del mínimo. Además, la demanda diaria presenta un ciclo cercano a un mes, compatible con el ritmo institucional de evaluaciones, y el comportamiento de uso persiste de manera marcada en Cálculo I después de que el PPA desaparece.

Dos resultados moderan una interpretación simplista. Primero, el desempeño de quienes tienen exactamente tres visitas es muy parecido al de quienes tienen cuatro o más, lo que sugiere que el contraste más fuerte está entre no uso/uso limitado y uso efectivo, no necesariamente en cada visita adicional. Segundo, la regularidad temporal de las visitas no predice desempeño de forma independiente una vez controlada la intensidad.

En términos científicos, el hallazgo central puede sintetizarse como
\[
\boxed{
\text{uso del CMAT}\;\longleftrightarrow\;\text{mejor desempeño}
\quad+\quad
\text{uso inicial}\;\longleftrightarrow\;\text{persistencia posterior},
}
\]
con una reserva esencial:
\[
\boxed{\text{asociación observacional}\neq\text{efecto causal identificado}.}
\]
La contribución potencial del trabajo consiste precisamente en mostrar la estructura temporal, académica y longitudinal de esa asociación con datos administrativos reales, mientras se explicitan los mecanismos rivales que la evidencia disponible todavía no puede separar.

% Esta sección se agrega de forma aditiva al borrador principal.
% No sustituye ni acorta las secciones anteriores.
\section{Experimentos adicionales: agrupación 1--2, licenciatura y ANOVA de Welch}
\label{sec:experimentos_adicionales}

\subsection{Convención del outcome y unidad de comparación}

En todos los análisis de esta sección se conserva exactamente la definición de desempeño continuo utilizada en el resto del proyecto. Para el estudiante $i$ perteneciente al salón $c$,
\[
Z_{ic}=\frac{Y_{ic}-\bar Y_c}{s_c},
\]
donde $Y_{ic}$ es la calificación final bajo la especificación primaria del outcome, $\bar Y_c$ y $s_c$ son la media y la desviación estándar del salón, y un \emph{salón} se define como un profesor impartiendo la misma materia en el mismo periodo académico. Como la cohorte primaria de esta sección contiene solamente el primer intento de \textit{Matemáticas Universitarias}, en la práctica $c$ equivale a profesor $\times$ periodo. Por construcción, $Z>0$ significa desempeño por encima del promedio del propio contexto de evaluación y $Z<0$ significa desempeño por debajo de ese promedio.

La unidad primaria sigue siendo un estudiante en su primer intento. Por tanto, cada estudiante contribuye una sola observación a estas comparaciones. Esta restricción elimina la dependencia mecánica que surgiría si un mismo estudiante aportara varios intentos de la misma materia.

\subsection{Por qué se agrupan una y dos visitas en la cohorte \texorpdfstring{$1$--$2$}{1--2}}

\subsubsection{Razón académica}

La razón institucional es directa: tanto una visita como dos visitas permanecen por debajo del umbral de tres visitas necesario para alcanzar el PPA. En consecuencia, $V=1$ y $V=2$ representan dos niveles de uso limitado que todavía no satisfacen el incentivo institucional. Esta similitud sustantiva justifica considerar su agrupación, pero por sí sola no demuestra que sus resultados académicos sean estadísticamente comparables.

\subsubsection{Prueba estadística de equivalencia}

No se utiliza la ausencia de significancia de una prueba de diferencia como evidencia de igualdad. En su lugar, se preespecifica una prueba de equivalencia TOST (\emph{two one-sided tests}) con margen práctico de $\pm0.20$ unidades $Z$. Si
\[
\Delta=\mu_{V=2}-\mu_{V=1},
\]
la hipótesis de no equivalencia se descompone en
\[
H_{01}:\Delta\leq -0.20,
\qquad
H_{02}:\Delta\geq 0.20.
\]
Se concluye equivalencia práctica de medias únicamente si ambas hipótesis se rechazan al nivel $\alpha=0.05$. El margen de $0.20Z$ debe interpretarse como una diferencia estandarizada pequeña elegida como \emph{smallest effect size of interest} para este borrador; no constituye una constante universal.

En los datos observados, el grupo con exactamente una visita contiene $n_1=517$ estudiantes y tiene media $\bar Z_1=0.222$, mientras que el grupo con exactamente dos visitas contiene $n_2=245$ y tiene media $\bar Z_2=0.209$. La diferencia estimada, orientada como $V=2$ menos $V=1$, es
\[
\widehat{\Delta}=-0.013Z,
\qquad IC_{95\%}=[-0.135,0.108].
\]
La prueba $t$ de Welch no detecta una diferencia de medias ($t=-0.216$, $\nu=511.6$, $p=0.8293$). Como comprobaciones distributivas, Mann--Whitney produce $p=0.6460$ y Brunner--Munzel $p=0.6423$; el delta de Cliff es -0.021, cercano a cero. La prueba de Brown--Forsythe para dispersión también es compatible con varianzas semejantes entre estos dos grupos ($p=0.6043$), aunque ninguna de las pruebas principales depende de suponer varianzas iguales.

La evidencia relevante para justificar la agrupación proviene del TOST: $p_{\mathrm{TOST}}=0.0013$. Por tanto, con el margen preespecificado $\pm0.20Z$, se rechaza que la diferencia verdadera de medias sea de magnitud igual o mayor a $0.20Z$ en cualquiera de las dos direcciones. En este sentido preciso, $V=1$ y $V=2$ son estadísticamente equivalentes en media para el objetivo operativo del estudio.

\begin{answerbox}
La cohorte \textbf{1--2 visitas} queda justificada por dos argumentos independientes: (i) ambas intensidades permanecen por debajo del umbral institucional de tres visitas y, por tanto, comparten una interpretación académica; y (ii) una prueba formal de equivalencia respalda que sus medias de desempeño difieren en menos de $0.20$ desviaciones estándar de salón. Esta conclusión se refiere a equivalencia práctica de la media, no a identidad completa de las distribuciones ni a equivalencia causal de los estudiantes que integran ambos grupos.
\end{answerbox}

\subsection{ANOVA de Welch entre las cuatro cohortes de visitas}

La comparación global utiliza las cohortes $0$, $1$--$2$, $3$ y $4+$. La hipótesis nula es
\[
H_0:\mu_0=\mu_{1-2}=\mu_3=\mu_{4+}.
\]
No se usa ANOVA clásico porque las varianzas empíricas de los grupos no son iguales. La prueba de Brown--Forsythe, centrada en la mediana, se utiliza como diagnóstico de heterocedasticidad. Dado que las varianzas y los tamaños de grupo son distintos, la prueba principal es el ANOVA de Welch, que reemplaza la varianza común por estimaciones específicas de cada grupo y ajusta los grados de libertad. Las comparaciones post hoc se realizan con Games--Howell, método compatible con varianzas y tamaños muestrales desiguales y que controla el error familiar de las comparaciones por pares.

El diagnóstico de Brown--Forsythe rechaza homogeneidad de varianzas ($W=20.735$, $p=2.29\times10^{-13}$), por lo que la elección de Welch es substantivamente pertinente. El ANOVA de Welch muestra diferencias globales claras entre las cuatro cohortes,
\[
F_W(3,546.4)=56.211,\qquad p=1.11\times10^{-31}.
\]

\begin{table}[H]
\centering
\caption{Desempeño continuo por cohorte de visitas para el ANOVA de Welch.}
\begin{tabular}{lrrrr}
\toprule
Grupo & $n$ & Media $Z$ & DE & Mediana $Z$ \\
\midrule
0 & 5393 & -0.060 & 1.017 & 0.315 \\
1-2 & 762 & 0.218 & 0.817 & 0.453 \\
3 & 170 & 0.325 & 0.724 & 0.543 \\
4+ & 302 & 0.342 & 0.712 & 0.535 \\
\bottomrule
\end{tabular}
\end{table}

Las comparaciones Games--Howell se muestran a continuación. Una diferencia negativa significa que la media del primer grupo es menor que la del segundo.

\begin{table}[H]
\centering
\caption{Comparaciones post hoc Games--Howell entre cohortes de visitas.}
\begin{tabular}{llrrr}
\toprule
Grupo A & Grupo B & $\Delta_{A-B}$ & $p_{GH}$ & IC95\% \\
\midrule
0 & 1-2 & -0.278 & $<0.001$ & [-0.362, -0.194] \\
0 & 3 & -0.385 & $<0.001$ & [-0.534, -0.237] \\
0 & 4+ & -0.402 & $<0.001$ & [-0.514, -0.291] \\
1-2 & 3 & -0.107 & 0.327 & [-0.269, 0.056] \\
1-2 & 4+ & -0.124 & 0.069 & [-0.254, 0.006] \\
3 & 4+ & -0.017 & 0.900 & [-0.195, 0.161] \\
\bottomrule
\end{tabular}
\end{table}

Los tres contrastes que involucran al grupo de cero visitas frente a $1$--$2$, $3$ o $4+$ son significativos después del ajuste Games--Howell. En cambio, no se detectan diferencias post hoc entre $1$--$2$ y $3$ ($p=0.327$), entre $1$--$2$ y $4+$ ($p=0.069$), ni entre $3$ y $4+$ ($p\geq0.90$). Este patrón es compatible con que una parte importante de la separación global ocurra entre no utilizar el servicio y utilizarlo, mientras que las diferencias adicionales entre niveles positivos de uso son menores e imprecisas.

Como el ANOVA de Welch es una comparación marginal y trata a los estudiantes como observaciones independientes, se añade una comprobación que reconoce la estructura compartida de salón. Un modelo lineal con efectos fijos profesor--periodo y matriz de covarianza robusta agrupada por el mismo salón produce un contraste Wald conjunto de los tres términos de cohorte de $W=153.193$ con $3$ restricciones y $p=5.39\times10^{-33}$ sobre 190 salones. La conclusión global, por tanto, no depende exclusivamente de ignorar la agrupación profesor--periodo.

\begin{figure}[H]
\centering
\includegraphics[width=0.72\textwidth]{figures/03_primary_outcome_by_visit_group}
\caption{Media de desempeño estandarizado por cohorte de visitas. El ANOVA de Welch evalúa la igualdad global de estas medias sin asumir varianzas iguales.}
\end{figure}

\subsection{Licenciatura, uso del CMAT y desempeño relativo al salón}

La licenciatura puede ser relevante por dos mecanismos distintos. Primero, estudiantes de programas distintos pueden tener distinta propensión a utilizar el CMAT. Segundo, aun después de expresar la calificación respecto al salón mediante $Z$, las distribuciones de desempeño pueden diferir entre licenciaturas. Si ambas asociaciones están presentes, la licenciatura es una covariable observada relevante para el análisis del vínculo entre asesorías y desempeño.

Para evitar inferencia dominada por programas con tamaños extremadamente pequeños, la comparación omnibus primaria conserva las licenciaturas con al menos 30 estudiantes con outcome continuo válido. Esta regla retiene 29 licenciaturas y 6,374 estudiantes, es decir, aproximadamente 96.2\% de la cohorte primaria. Las licenciaturas de menor tamaño permanecen en los archivos descriptivos completos, pero no se utilizan en el ANOVA principal.

\subsubsection{Licenciatura y desempeño estandarizado}

La prueba de Brown--Forsythe vuelve a mostrar heterogeneidad de varianzas entre licenciaturas ($W=2.091$, $p=0.0006$). El ANOVA de Welch rechaza igualdad global de medias de $Z$,
\[
F_W(28,1239.4)=4.308,\qquad p=9.89\times10^{-13}.
\]
Por tanto, incluso después de definir el desempeño relativamente al profesor y periodo del estudiante, la posición promedio dentro del salón no es idéntica entre licenciaturas.

Como sensibilidad frente a dependencia residual dentro de salón, un modelo con efectos fijos profesor--periodo y errores estándar agrupados por salón realiza un contraste Wald conjunto sobre 28 términos de licenciatura. El resultado es $W=165.211$, $p=2.11\times10^{-21}$. La incorporación de licenciatura aumenta $R^2$ de 0.0016 a 0.0247, un incremento de 0.0232. La magnitud explicada es modesta, aunque estadísticamente detectable en esta muestra grande; por ello la licenciatura debe tratarse como una fuente de heterogeneidad real pero no como explicación dominante del desempeño individual.

\begin{figure}[H]
\centering
\includegraphics[width=0.86\textwidth]{figures/12_career_classroom_z_means}
\caption{Media de desempeño $Z$ respecto al salón por licenciatura. Se muestran únicamente licenciaturas con $n\geq30$ en la figura inferencial.}
\end{figure}

\subsubsection{Licenciatura y probabilidad de utilizar el CMAT}

Para una primera prueba con supuestos clásicos manejables, se reduce la exposición a ``al menos una visita'' frente a ninguna visita. La tabla licenciatura $\times$ uso tiene conteos esperados mayores que cinco en todas sus celdas. El contraste de Pearson obtiene $\chi^2(28)=92.618$, $p=7.69\times10^{-9}$, con $V$ de Cramér corregido por sesgo igual a 0.101. La asociación es estadísticamente clara pero de magnitud pequeña.

Al utilizar las cuatro cohortes $0$, $1$--$2$, $3$ y $4+$, 20.7\% de las celdas esperadas quedan por debajo de cinco, de modo que la aproximación asintótica de $\chi^2$ ya no es ideal. Por esta razón se conserva el estadístico de Pearson como descriptivo, pero el $p$ principal se obtiene por permutación de las etiquetas de cohorte manteniendo los márgenes observados. Con 3000 permutaciones se obtiene $p_{perm}=0.0003$ y $V$ de Cramér corregido igual a 0.056.

\begin{figure}[H]
\centering
\includegraphics[width=0.86\textwidth]{figures/14_career_any_visit_rates}
\caption{Proporción de estudiantes con al menos una visita al CMAT por licenciatura, restringida a programas con $n\geq30$.}
\end{figure}

Estos dos resultados justifican mantener la licenciatura como covariable observada en los modelos ajustados: está asociada tanto con la exposición (uso del CMAT) como con el outcome (posición estandarizada dentro del salón). Esto no elimina autoselección no observada, pero evita atribuir al uso del CMAT diferencias que sean explicables por composición de licenciaturas.

\begin{longtable}{lrrrrr}
\caption{Resumen por licenciatura para programas con al menos 30 estudiantes.}\label{tab:career_summary_extended}\\
\toprule
Lic. & $n$ & Media $Z$ & DE $Z$ & Uso $\geq1$ & Media visitas \\
\midrule
\endfirsthead
\toprule
Lic. & $n$ & Media $Z$ & DE $Z$ & Uso $\geq1$ & Media visitas \\
\midrule
\endhead
LBM & 199 & 0.246 & 0.885 & 21.1\% & 0.528 \\
LNM & 249 & 0.190 & 0.933 & 17.3\% & 0.494 \\
LII & 361 & 0.183 & 0.805 & 19.1\% & 0.640 \\
LBC & 118 & 0.168 & 0.816 & 23.7\% & 0.644 \\
LDS & 87 & 0.165 & 0.985 & 26.4\% & 1.092 \\
LMT & 320 & 0.146 & 0.899 & 14.4\% & 0.347 \\
LCS & 86 & 0.134 & 1.005 & 24.4\% & 0.558 \\
LAT & 368 & 0.114 & 0.964 & 30.4\% & 0.918 \\
LEC & 174 & 0.098 & 0.934 & 16.1\% & 0.580 \\
LNI & 550 & 0.089 & 0.885 & 19.8\% & 0.582 \\
LBN & 281 & 0.060 & 0.988 & 16.0\% & 0.555 \\
LQI & 151 & 0.051 & 1.038 & 22.5\% & 0.801 \\
LNA & 53 & 0.050 & 1.005 & 22.6\% & 0.509 \\
LIN & 64 & 0.043 & 0.966 & 20.3\% & 0.406 \\
LFA & 138 & 0.029 & 0.948 & 24.6\% & 0.659 \\
LIC & 222 & 0.021 & 0.995 & 14.4\% & 0.432 \\
LIM & 201 & -0.004 & 1.049 & 12.4\% & 0.393 \\
LIS & 271 & -0.005 & 1.054 & 16.2\% & 0.443 \\
LRT & 96 & -0.010 & 1.049 & 16.7\% & 1.000 \\
LBI & 152 & -0.013 & 1.002 & 18.4\% & 0.414 \\
LAR & 445 & -0.024 & 0.960 & 9.7\% & 0.234 \\
LFP & 343 & -0.032 & 1.009 & 19.8\% & 0.551 \\
LCN & 160 & -0.063 & 0.907 & 18.8\% & 0.438 \\
LFB & 131 & -0.075 & 0.997 & 23.7\% & 0.641 \\
LMK & 350 & -0.080 & 1.002 & 20.0\% & 0.597 \\
LHR & 178 & -0.131 & 1.075 & 15.2\% & 0.534 \\
LIY & 34 & -0.202 & 1.251 & 14.7\% & 0.147 \\
LAD & 562 & -0.266 & 0.999 & 20.8\% & 0.523 \\
LPM & 30 & -0.631 & 1.383 & 23.3\% & 0.600 \\
\bottomrule
\end{longtable}


\subsection{Heterogeneidad exploratoria de la asociación por licenciatura}

La existencia simultánea de diferencias de uso y de desempeño entre licenciaturas sugiere una pregunta adicional: ¿la asociación entre cohorte de visitas y $Z$ es numéricamente homogénea entre programas? Este análisis se considera exploratorio y no se utiliza para redefinir el estimando primario.

Para evitar interacciones sustentadas en celdas casi vacías, se conservaron únicamente licenciaturas con al menos 100 estudiantes y al menos cinco observaciones en cada cohorte $0$, $1$--$2$, $3$ y $4+$. Quedan 19 programas y $N=5{,}481$. Se estima un modelo sobre el outcome ya estandarizado por salón,
\[
Z_i = \alpha + \beta_{G(i)} + \gamma_{L(i)} + (\beta\gamma)_{G(i),L(i)} + \varepsilon_i,
\]
con matriz de covarianza robusta agrupada por profesor--periodo. El contraste Wald conjunto de los 54 términos de interacción produce $W=443.96$ y $p=1.12\times10^{-62}$.

\begin{interpretbox}
La asociación entre utilización y desempeño no parece numéricamente idéntica entre todas las licenciaturas. Sin embargo, este resultado es de alta dimensión y debe tratarse como evidencia de heterogeneidad a investigar, no como prueba de que una licenciatura específica ``se beneficia más'' del CMAT. La versión final deberá preespecificar qué contrastes por programa tienen justificación académica antes de interpretar efectos particulares.
\end{interpretbox}

\subsection{Supuestos, alcance y precauciones de estas pruebas}

\begin{enumerate}[label=\arabic*.]
\item \textbf{Independencia entre estudiantes.} La cohorte contiene un único primer intento por alumno, por lo que no existen repeticiones del mismo individuo dentro de estas comparaciones. Sin embargo, alumnos del mismo profesor--periodo comparten contexto de evaluación. Por ello, los ANOVA de Welch se interpretan como pruebas marginales y se acompañan de contrastes Wald con covarianza agrupada por salón.
\item \textbf{Heterocedasticidad.} No se asume igualdad de varianzas. Brown--Forsythe se reporta como diagnóstico; Welch y Games--Howell son precisamente las versiones seleccionadas para trabajar bajo dispersión desigual y tamaños de grupo muy diferentes.
\item \textbf{Forma de las distribuciones.} Welch requiere medias con distribución muestral suficientemente regular y varianzas finitas, no normalidad exacta de cada observación. Los tamaños grandes de las cohortes principales favorecen la aproximación asintótica. En la comparación $1$ frente a $2$ se añaden Mann--Whitney y Brunner--Munzel como comprobaciones menos dependientes de la forma paramétrica.
\item \textbf{Equivalencia.} El TOST no afirma que los grupos $V=1$ y $V=2$ sean idénticos. Afirma que la diferencia de medias es compatible con una banda de equivalencia preespecificada de $\pm0.20Z$. Cambiar ese margen cambiaría la pregunta científica.
\item \textbf{Comparaciones múltiples.} Después del ANOVA de Welch, los pares se evalúan con Games--Howell; no se realizan seis pruebas $t$ independientes sin corrección.
\item \textbf{Licenciaturas pequeñas.} La inferencia entre programas se restringe a $n\geq30$. Los programas pequeños se conservan descriptivamente en los CSV, pero sus estimaciones no deben interpretarse con la misma estabilidad.
\item \textbf{Chi-cuadrada y celdas esperadas.} Para cualquier visita vs. ninguna, las celdas esperadas satisfacen el criterio clásico de al menos cinco. Para la tabla de cuatro cohortes, donde ese supuesto se incumple en una fracción de celdas, el $p$ principal se obtiene por permutación.
\item \textbf{No causalidad.} Ninguna de estas pruebas convierte la asistencia voluntaria en tratamiento aleatorio. Motivación, preparación matemática, percepción de dificultad y organización personal permanecen potencialmente no observadas.
\end{enumerate}

\paragraph{Síntesis.}
Los experimentos adicionales fortalecen tres decisiones del borrador. Primero, la categoría $1$--$2$ tiene una justificación institucional y una justificación estadística formal de equivalencia. Segundo, las cuatro cohortes no tienen la misma media global, aun permitiendo varianzas distintas; la diferencia más clara se concentra entre cero visitas y cualquier nivel positivo de uso. Tercero, la licenciatura está asociada tanto con la utilización del CMAT como con la posición académica relativa dentro del salón, por lo que debe conservarse como covariable de ajuste y como dimensión de heterogeneidad en análisis posteriores.

% -----------------------------------------------------------------------------
% Extensiones añadidas en v7. Se conservan intactos los análisis anteriores.
% -----------------------------------------------------------------------------
\subsection{Índice de rendimiento por número exacto de asesorías: 1 a 12 visitas}
\label{subsec:indice_exacto_1_12}

\subsubsection{Motivación y definición del índice}

El análisis anterior utiliza cohortes amplias ($0$, $1$--$2$, $3$ y $4+$) porque son apropiadas para responder preguntas ligadas al umbral PPA y evitan celdas excesivamente pequeñas. Sin embargo, esa agrupación oculta la forma detallada de la relación entre intensidad de uso y desempeño. Como extensión descriptiva se reconstruyen cohortes por número \emph{exacto} de visitas al CMAT durante el mismo semestre de Matemáticas Universitarias,
\[
V_i\in\{1,2,\ldots,12\}.
\]
No se incluyen en esta figura los estudiantes con cero visitas porque el objetivo de este índice es describir la variación de desempeño entre usuarios con distintas intensidades observadas. Los estudiantes con más de 12 visitas permanecen en la base y en los análisis principales; simplemente quedan fuera de esta visualización de dosis exacta para evitar una cola de conteos extremadamente dispersa.

El punto de partida sigue siendo el outcome principal estandarizado dentro del salón,
\[
Z_{ic}=\frac{Y_{ic}-\bar Y_c}{s_c},
\qquad
c=\text{profesor}\times\text{misma materia}\times\text{mismo periodo}.
\]
Como el análisis por licenciatura mostró diferencias estadísticamente detectables en la media de $Z$ entre programas, se añade una segunda normalización para expresar el desempeño de cada estudiante respecto a su propia licenciatura. Para la licenciatura $\ell$ se define
\[
R_i=\frac{Z_i-\bar Z_{\ell(i)}}{s_{\ell(i)}}.
\]
Así, $R_i=0.30$ significa que el estudiante se encuentra aproximadamente 0.30 desviaciones estándar por encima del promedio de estudiantes de su misma licenciatura, \emph{después} de haber expresado primero su calificación respecto al salón. Esta segunda estandarización no sustituye al $Z$ de salón; funciona como un índice complementario para evitar que las diferencias de composición entre licenciaturas dominen la comparación de conteos exactos de visitas.

Para estabilizar la media y desviación estándar utilizadas en esta segunda normalización, se conservan las licenciaturas con al menos 30 estudiantes con outcome válido. Son las mismas 29 licenciaturas utilizadas en el análisis inferencial anterior. El índice por número exacto de visitas es entonces
\[
I_v=E(R_i\mid V_i=v),\qquad v=1,\ldots,12.
\]
Se reporta una versión ponderada por estudiante y, como sensibilidad descriptiva, una versión balanceada por licenciatura en la que cada programa con al menos dos estudiantes dentro de la celda de visitas recibe el mismo peso. La primera responde a ``¿cómo se desempeña el estudiante promedio observado con $v$ visitas respecto a su licenciatura?''; la segunda reduce el dominio mecánico de las licenciaturas más grandes.

\subsubsection{Resultados del índice exacto}

\begin{table}[H]
\centering
\caption{Índice de rendimiento por número exacto de visitas. $Z$ se define respecto al salón e $I_v$ respecto a la licenciatura. Los IC95\% de $I_v$ son intervalos $t$ descriptivos.}
\label{tab:exact_visit_index}
\begin{tabular}{rrrrr}
\toprule
Visitas & $n$ & Media $Z$ salón & $I_v$ licenciatura & IC95\% de $I_v$ \\
\midrule
1  & 500 & 0.233 & 0.220 & [0.145, 0.295] \\
2  & 239 & 0.229 & 0.217 & [0.119, 0.314] \\
3  & 168 & 0.322 & 0.304 & [0.200, 0.409] \\
4  & 94  & 0.261 & 0.241 & [0.091, 0.391] \\
5  & 53  & 0.343 & 0.354 & [0.150, 0.559] \\
6  & 48  & 0.311 & 0.311 & [0.059, 0.563] \\
7  & 25  & 0.184 & 0.165 & [-0.189, 0.520] \\
8  & 19  & 0.322 & 0.284 & [-0.085, 0.652] \\
9  & 11  & 0.507 & 0.546 & [0.234, 0.858] \\
10 & 13  & 0.442 & 0.453 & [0.198, 0.708] \\
11 & 2   & 0.669 & 0.686 & [-0.471, 1.844] \\
12 & 4   & 0.815 & 0.762 & [0.207, 1.318] \\
\bottomrule
\end{tabular}
\end{table}

\begin{figure}[H]
\centering
\includegraphics[width=0.88\textwidth]{figures/11_exact_visit_career_relative_index}
\caption{Índice de rendimiento relativo a la licenciatura según el número exacto de visitas al CMAT durante el semestre. Las barras son IC95\% $t$ descriptivos; los tamaños muestrales aparecen junto a cada punto.}
\label{fig:exact_visit_index}
\end{figure}

Los primeros seis conteos presentan índices promedio positivos de aproximadamente $0.22$ a $0.35$ desviaciones estándar relativas a la licenciatura. A partir de siete visitas la incertidumbre aumenta rápidamente debido a la reducción del tamaño muestral. En particular, las cohortes de 8 a 12 visitas contienen 19, 11, 13, 2 y 4 estudiantes, respectivamente. Los valores altos observados en 9--12 visitas no deben interpretarse como evidencia de un aumento sistemático del rendimiento con cada asesoría adicional: esos puntos tienen gran incertidumbre y pueden ser sensibles a pocos estudiantes.

La ausencia de una relación monotónica clara se confirma con dos análisis de tendencia preespecificados como exploratorios. Utilizando todos los estudiantes con 1--12 visitas y el índice individual $R_i$, la correlación de Spearman es
\[
\rho=0.033,\qquad p=0.255,
\]
y una regresión lineal de $R_i$ sobre el número exacto de visitas, con errores estándar agrupados por salón, estima
\[
\widehat\beta=0.0225\;\text{DE por visita},
\qquad IC_{95\%}=[-0.0005,0.0455],
\qquad p=0.055.
\]
Como sensibilidad frente a la cola de conteos escasos, se repite el análisis utilizando únicamente las cohortes exactas con al menos 20 observaciones. Esta regla conserva 1--7 visitas ($N=1{,}127$). La correlación de Spearman es entonces $\rho=0.022$ ($p=0.461$) y la pendiente lineal agrupada es $0.0151$ DE por visita (IC95\% $[-0.0225,0.0527]$, $p=0.431$).

\begin{answerbox}
El índice exacto muestra que los usuarios del CMAT se ubican, en promedio, por encima de la media de su propia licenciatura en varios niveles de intensidad. Sin embargo, los datos no muestran una relación monotónica o lineal estadísticamente clara entre 1 y 12 visitas. Esto es consistente con el resultado previo de una posible meseta después de niveles bajos/moderados de uso y evita interpretar el conteo de asesorías como una ``dosis'' con retorno marginal constante.
\end{answerbox}

\subsubsection{Supuestos y límites del índice}

\begin{enumerate}[label=\arabic*.]
\item \textbf{Dos niveles de estandarización.} El primer nivel ($Z$) controla escala y dispersión del contexto profesor--materia--periodo; el segundo ($R$) expresa la posición de ese $Z$ dentro de la licenciatura. Por ello $R$ no es una calificación original ni una medida causal del efecto del CMAT.
\item \textbf{Comparación con la licenciatura.} La media y desviación estándar de cada licenciatura se calculan usando la cohorte primaria completa de esa licenciatura, no solamente a quienes acudieron al CMAT. Esto permite interpretar cero como el promedio del programa en la métrica $Z$ de salón.
\item \textbf{Tamaños pequeños.} Los intervalos de la tabla son $t$ de Student para la media porque la aproximación normal de 1.96 es especialmente inadecuada en las cohortes más pequeñas. Aun así, con $n=2$ o $n=4$ los intervalos son necesariamente inestables y se consideran descriptivos.
\item \textbf{Dependencia por salón.} Los intervalos por punto son descriptivos y no corrigen agrupación. La prueba de tendencia complementaria sí utiliza errores estándar robustos agrupados por profesor--periodo.
\item \textbf{No linealidad.} La pendiente lineal no se interpreta como modelo estructural de ``ganancia por visita''. Se reporta junto con Spearman y con una sensibilidad que excluye celdas escasas precisamente porque la forma de la relación puede ser no lineal.
\item \textbf{Autoselección.} Estudiantes que acuden muchas veces pueden diferir en necesidad, motivación, disponibilidad, organización o preparación previa. Por ello $I_v$ es un índice de asociación entre intensidad observada y desempeño relativo, no una curva causal dosis--respuesta.
\end{enumerate}

\subsection{Persistencia individual de uso de Matemáticas Universitarias a Cálculo I}
\label{subsec:mu_calc_2x2}

El análisis longitudinal anterior mostró un gradiente de utilización posterior según las cohortes $0$, $1$--$2$, $3$ y $4+$. Aquí se formula una pregunta más elemental y directamente interpretable: entre estudiantes que progresan posteriormente a Cálculo I y para quienes existe cobertura de asesorías en ambos periodos, ¿haber utilizado alguna vez el CMAT durante Matemáticas Universitarias está asociado con volver a utilizarlo durante Cálculo I?

Se definen dos indicadores binarios,
\[
M_i=1(V^{MU}_i>0),
\qquad
C_i=1(V^{Calc}_i>0).
\]
La cohorte longitudinal analizable contiene $N=4{,}211$ estudiantes y produce las cuatro combinaciones de la Tabla~\ref{tab:mu_calc_2x2}.

\begin{table}[H]
\centering
\caption{Combinaciones individuales de cualquier utilización del CMAT en Matemáticas Universitarias y Cálculo I.}
\label{tab:mu_calc_2x2}
\begin{tabular}{lrr}
\toprule
Patrón & $n$ & Proporción de la cohorte \\
\midrule
No MU / no Cálculo & 2721 & 64.62\% \\
No MU / sí Cálculo & 553 & 13.13\% \\
Sí MU / no Cálculo & 492 & 11.68\% \\
Sí MU / sí Cálculo & 445 & 10.57\% \\
\bottomrule
\end{tabular}
\end{table}

La probabilidad observada de utilizar el CMAT en Cálculo entre quienes ya lo habían utilizado en MU es
\[
P(C_i=1\mid M_i=1)=0.4749,
\qquad IC_{95\%}^{Wilson}=[0.4431,0.5069],
\]
mientras que entre quienes no lo habían utilizado es
\[
P(C_i=1\mid M_i=0)=0.1689,
\qquad IC_{95\%}^{Wilson}=[0.1565,0.1821].
\]
La diferencia absoluta de riesgos es de 0.306, es decir, 30.6 puntos porcentuales, con IC95\% de Newcombe--Wilson $[0.272,0.340]$. El riesgo relativo es
\[
RR=2.812,
\qquad IC_{95\%}=[2.540,3.112],
\]
y la razón de momios es
\[
OR=4.450,
\qquad IC_{95\%}=[3.802,5.209].
\]
La prueba de Pearson para la tabla $2\times2$ obtiene $\chi^2(1)=377.26$ con $p=4.91\times10^{-84}$; el coeficiente $\phi=0.299$. Todas las frecuencias esperadas son grandes (mínimo $=222.1$), por lo que la aproximación asintótica es apropiada. Como comprobación exacta, Fisher produce $p=3.58\times10^{-76}$ y la misma razón de momios puntual.

\begin{figure}[H]
\centering
\includegraphics[width=0.78\textwidth]{figures/13_any_visit_mu_to_calculus_transition}
\caption{Probabilidad observada de utilizar el CMAT durante Cálculo I según utilización previa durante Matemáticas Universitarias.}
\label{fig:mu_calc_transition}
\end{figure}

Hay una distinción importante entre \emph{persistencia individual} y \emph{cambio en la prevalencia total}. En esta misma cohorte, 22.25\% utilizó el CMAT durante MU y 23.70\% lo utilizó durante Cálculo. La prueba de McNemar, que compara específicamente los dos tipos de estudiantes discordantes (553 que comienzan a utilizarlo en Cálculo frente a 492 que dejan de utilizarlo), produce $\chi^2_{McN}=3.445$ y $p=0.0634$. Por tanto, no se detecta al 5\% un cambio marginal claro en la proporción total de usuarios entre ambos cursos. Esto no contradice la asociación longitudinal: la prevalencia agregada puede permanecer aproximadamente estable y, al mismo tiempo, el uso concentrarse fuertemente en estudiantes que ya tenían experiencia previa con el servicio.

\begin{answerbox}
Entre estudiantes que llegan posteriormente a Cálculo I, quienes utilizaron el CMAT durante Matemáticas Universitarias tienen aproximadamente 2.8 veces la probabilidad observada de utilizarlo de nuevo. El hallazgo describe una fuerte persistencia individual del patrón de \emph{help-seeking}. Al mismo tiempo, McNemar muestra que la proporción total de usuarios no cambia de forma estadísticamente clara entre ambos cursos; la señal principal está en \emph{quién} utiliza el servicio, no simplemente en cuántos lo utilizan.
\end{answerbox}

\subsubsection{Supuestos y alcance de la transición MU--Cálculo}

\begin{enumerate}[label=\arabic*.]
\item \textbf{Condicionamiento por progresión.} El análisis incluye únicamente estudiantes que posteriormente cursan Cálculo I y tienen cobertura de asesorías. No representa a todos los estudiantes que comenzaron MU; progresar a Cálculo depende de la trayectoria académica previa.
\item \textbf{Asociación, no causalidad.} $RR$, $OR$ y la diferencia de riesgos cuantifican persistencia observada. No demuestran que utilizar el CMAT en MU \emph{cause} la utilización posterior. Una propensión estable a pedir ayuda, la necesidad académica u otras características individuales pueden explicar parte de la asociación.
\item \textbf{Prueba de Pearson/Fisher.} Estas pruebas evalúan asociación entre uso previo y posterior. La unidad es el estudiante y cada persona aparece una sola vez en la tabla.
\item \textbf{McNemar responde otra pregunta.} McNemar aprovecha el carácter pareado para evaluar si cambia la prevalencia marginal de uso de MU a Cálculo. No es una prueba de la asociación $M_i\leftrightarrow C_i$, por lo que su $p=0.0634$ no invalida el $RR$ ni el $OR$.
\item \textbf{Cobertura temporal.} Un cero significa ausencia de visitas únicamente en periodos para los que existe registro efectivo de asesorías; los periodos sin cobertura no se codifican como no uso.
\end{enumerate}

\paragraph{Síntesis de las extensiones.}
El índice exacto 1--12 y la tabla longitudinal $2\times2$ añaden dos piezas complementarias. Primero, la ventaja relativa de los usuarios respecto a su propia licenciatura no se traduce en una pendiente monotónica convincente por cada visita adicional: la intensidad exacta debe interpretarse con cautela, especialmente en la cola de 8--12 visitas. Segundo, el comportamiento de búsqueda de apoyo es fuertemente persistente a nivel individual entre MU y Cálculo, aun cuando la prevalencia total de usuarios de ambos cursos sea similar. En conjunto, estos resultados favorecen una descripción del CMAT en términos de adopción, utilización recurrente y heterogeneidad estudiantil, más que una narrativa simple de ``más visitas implican linealmente mayor desempeño''.

\clearpage
\section{Nueva arquitectura longitudinal: PPA1, adopción, persistencia y trayectoria hacia Cálculo I}
\label{sec:ppa_progression_v8}

Esta sección se agrega como un capítulo de trabajo independiente. No sustituye ni corrige silenciosamente los análisis anteriores del borrador: los conserva como antecedentes exploratorios y formula una arquitectura longitudinal más estrecha a partir de información administrativa adicional sobre el funcionamiento del CMAT, el Programa de Primer Año (PPA) y las revalidaciones entre licenciaturas. La intención es acumular evidencia y decisiones de diseño antes de reorganizar el manuscrito definitivo para publicación.

\subsection{Contexto institucional del CMAT y del PPA1}

De acuerdo con la descripción institucional proporcionada para este proyecto, el Centro de Aprendizaje de Matemáticas (CMAT) atiende a estudiantes de licenciatura desde matemáticas básicas hasta cursos avanzados y ofrece apoyo en matemáticas y estadística. El portal institucional describe además diagnóstico, materiales de refuerzo, asesorías, talleres y conferencias, y señala como cursos atendidos principalmente Precálculo, Cálculo I, Cálculo II, Matemáticas I, Matemáticas II, distintas asignaturas de Estadística y Matemáticas Universitarias.\footnote{Centro de Aprendizaje de Matemáticas, UDLAP: \url{https://www.udlap.mx/centrodeMatematicas/conocenos.aspx}. La página pública actual también indica que actividades del CMAT pueden otorgar puntos PPA1. El presente estudio no utiliza como variable el número total de puntos otorgados por esa actividad; utiliza el umbral operacional de tres registros de visita que forma parte del diseño administrativo analizado.}

Cada asesoría genera un registro mediante Google Forms. El registro es capturado por estudiantes que realizan servicio social o servicio becario e incluye, entre otros campos, identificador del estudiante, licenciatura declarada, profesor consultado, materia y tema. Para este estudio se asume que dichos registros representan visitas reales, reconociendo que se trata de captura humana y que no existe un mecanismo independiente para garantizar ausencia absoluta de duplicados. Un mismo estudiante puede registrar más de una visita en el mismo día: los profesores trabajan por bloques horarios y un estudiante puede continuar sus dudas con otro profesor cuando termina el horario del primero. Por ello, varias visitas en una misma fecha se interpretan como \emph{concentración intradía}, no como evidencia automática de manipulación del sistema.

Todos los asesores del CMAT son profesores y cualquier estudiante vigente puede acudir libremente. El centro no se limita a MU: un estudiante avanzado podría solicitar apoyo para Cálculo, Ecuaciones Diferenciales, Probabilidad, Procesos Estocásticos u otra asignatura matemática/estadística mientras continúe inscrito. La restricción a Matemáticas Universitarias y Cálculo I es, por tanto, una decisión del diseño científico de este estudio y no una restricción de operación del CMAT.

Administrativamente, el PPA funciona como un requisito institucional de participación del primer año. Se distingue PPA1, asociado al primer semestre, y PPA2, asociado al segundo semestre. El estudiante acredita PPA1 acumulando puntos mediante distintas actividades universitarias, que pueden incluir organizaciones estudiantiles, eventos y actividades del CMAT. PPA1 no es parte de la calificación de Matemáticas Universitarias. Aunque teóricamente un estudiante podría completar obligaciones PPA1 de manera tardía, el análisis adopta explícitamente la siguiente suposición de trabajo:

\begin{cautionbox}[title=Suposición administrativa central del capítulo]
Se trata a Matemáticas Universitarias como el contexto típico de primer semestre/PPA1 y se asume que los estudiantes cumplen PPA1 durante ese primer semestre. En consecuencia, el umbral de tres registros de CMAT durante MU se interpreta como un umbral de utilización potencialmente influido por PPA1. Cálculo I se trata como el curso posterior en el que ese mismo incentivo PPA1 del CMAT ya no opera. Esta es una suposición del diseño y no una variable individual observada: no se conoce la motivación psicológica de cada estudiante ni se observa directamente qué punto PPA buscaba obtener en cada visita.
\end{cautionbox}

Esta precisión cambia la pregunta. No se intenta etiquetar a un estudiante como ``motivado por PPA'' o ``motivado académicamente''. Se observan patrones compatibles con diferentes conductas:
\[
\text{incentivo PPA1}\longrightarrow\text{adopción del CMAT}\longrightarrow
\text{continuación o abandono cuando desaparece el mismo incentivo}.
\]

\subsection{Dos fuentes de licenciatura y el problema de las revalidaciones}

Existen dos fuentes distintas para la licenciatura. La base académica oficial contiene un código de carrera asociado a cada registro de materia; el Google Form de asesorías contiene la licenciatura declarada por el estudiante en el momento de la visita. Estas variables no deben considerarse equivalentes de forma automática.

Un estudiante puede acumular más de una licenciatura. Además, cuando una materia previamente aprobada se revalida para un segundo programa, el sistema académico puede crear una fila nueva en un periodo posterior con la misma calificación histórica, aun cuando el estudiante no haya vuelto a cursar la materia. De esta manera,
\[
\boxed{\text{fila académica}\neq\text{nuevo intento real de la materia}.}
\]

El audit administrativo del archivo original identifica 10,080 estudiantes con una sola licenciatura oficial registrada, 448 con dos y 3 con tres. También identifica 537 filas posteriores a una aprobación previa del mismo curso, correspondientes a 263 estudiantes, que cumplen la regla conservadora de \emph{likely post-pass revalidation}; las 537 reproducen la misma calificación del primer aprobado observado. Se detectan además 405 filas redundantes dentro del mismo estudiante--materia--periodo con múltiples carreras, correspondientes a 231 estudiantes. Finalmente, 1,048 filas del extracto original carecen de profesor registrado. Esta última característica es coherente con el problema de revalidación: una acreditación administrativa puede conservar la calificación pero no representar un salón realmente cursado.

El código no borra estas observaciones del audit. Las clasifica. Para construir resultados académicos, una observación posterior al primer aprobado del mismo curso se marca como probable revalidación y no se usa como nuevo intento. Para construir el $Z$ de salón, se utilizan únicamente intentos reales candidatos con calificación numérica y profesor observado.

La licenciatura oficial de los modelos corresponde al registro del \emph{intento académico real} de MU o Cálculo, no a una fila posterior de revalidación. La licenciatura declarada en Google Forms se conserva como variable descriptiva contemporánea a la asesoría. No puede utilizarse como ajuste basal para comparar usuarios con no usuarios, porque un estudiante que nunca acudió al CMAT no tiene, por definición, un formulario de asesoría del cual recuperar esa licenciatura.

\subsection{Cohorte PPA1--Cálculo y definición estricta del Z-score}

Para las preguntas de este capítulo se adopta una cohorte más estrecha que la cohorte general del borrador. El estudiante debe:

\begin{enumerate}[label=\arabic*.,leftmargin=*]
\item tener un primer intento real observado de Matemáticas Universitarias;
\item aprobar ese primer intento con calificación numérica $Y_{MU}\geq7.5$;
\item cursar posteriormente un primer intento real de Cálculo I;
\item disponer de calificación numérica final en Cálculo I;
\item tener cobertura efectiva del archivo de asesorías en ambos periodos;
\item disponer de un salón con al menos cinco calificaciones numéricas y desviación estándar positiva en ambas materias;
\item para la especificación primaria de PPA1, avanzar a Cálculo en el siguiente periodo regular Primavera/Otoño.
\end{enumerate}

El flujo es:

\begin{table}[H]
\centering
\caption{Construcción de la cohorte longitudinal PPA1--Cálculo.}
\label{tab:ppa_v8_flow}
\begin{tabularx}{0.92\textwidth}{@{}Xr@{}}
\toprule
Etapa & $N$ \\
\midrule
Primer intento real candidato de MU & 7,777 \\
Primer intento de MU con aprobación numérica $\geq7.5$ & 5,769 \\
Existe un primer Cálculo I posterior & 4,556 \\
Cálculo I posterior con calificación numérica & 4,055 \\
Pares con cobertura CMAT y $Z$ válido en ambas materias & 3,389 \\
\textbf{Cohorte primaria: Cálculo en el siguiente periodo regular} & \textbf{3,241} \\
\bottomrule
\end{tabularx}
\end{table}

El análisis conserva además los 3,389 pares como sensibilidad que permite retrasos entre MU y Cálculo. Los resultados de persistencia cambian poco al permitir esos retrasos.

En este capítulo el $Z$ se calcula literalmente respecto al salón en el sentido administrativo acordado:
\[
\boxed{c=\text{profesor}\times\text{misma materia}\times\text{mismo periodo}.}
\]
Por tanto,
\[
Z_{i,MU}=\frac{Y_{i,MU}-\bar Y_{c,MU}}{s_{c,MU}},
\qquad
Z_{i,Calc}=\frac{Y_{i,Calc}-\bar Y_{c,Calc}}{s_{c,Calc}}.
\]
La media y desviación estándar de referencia se calculan usando las calificaciones numéricas de los intentos reales candidatos del salón, no únicamente los estudiantes que terminaron entrando a la cohorte longitudinal. Así, $Z$ representa posición relativa dentro del salón realmente observado y no dentro de una submuestra seleccionada posteriormente.

También se define
\[
\Delta Z_i=Z_{i,Calc}-Z_{i,MU},
\]
que describe un cambio de \emph{posición académica relativa}. No es un cambio de puntos de calificación y tampoco es, por sí mismo, un efecto causal del CMAT.

\subsection{Perfiles observables: uso limitado al contexto PPA y uso persistente}

Durante MU se utilizan las categorías ya establecidas:
\[
G_i\in\{0,\;1\text{--}2,\;3,\;4+\},
\]
y se define
\[
C_i=1(V_{i,Calc}>0).
\]
El cruce genera ocho perfiles observables. Las expresiones siguientes describen conductas registradas; no son diagnósticos de motivación interna.

\begin{table}[H]
\centering
\caption{Perfiles de utilización en la cohorte primaria PPA1--Cálculo ($N=3{,}241$).}
\label{tab:ppa_v8_profiles}
\begin{tabular}{@{}llr@{}}
\toprule
Uso en MU & Uso posterior en Cálculo & $n$ \\
\midrule
0 & No & 2,044 \\
0 & Sí: adopción tardía & 426 \\
1--2 & No & 267 \\
1--2 & Sí & 205 \\
Exactamente 3 & No: patrón compatible con uso limitado al umbral & 47 \\
Exactamente 3 & Sí: alcanza umbral y continúa después & 56 \\
4+ & No: uso por encima del umbral sin continuidad posterior & 65 \\
4+ & Sí: uso persistente por encima del umbral & 131 \\
\bottomrule
\end{tabular}
\end{table}

El grupo ``exactamente 3 y no Cálculo'' es el patrón más cercano a la hipótesis de utilización limitada al requisito, pero incluso aquí no se puede afirmar que el estudiante haya acudido \emph{solo} por los puntos. Del mismo modo, $4+$ no equivale a motivación intrínseca. Los nombres apropiados son \emph{threshold-limited pattern}, \emph{late adoption} y \emph{persistent use}.

\subsection{¿La utilización persiste cuando desaparece el mismo incentivo PPA1?}

\begin{questionbox}
Entre estudiantes que aprobaron MU y cursaron Cálculo I en el siguiente periodo regular, ¿la probabilidad de volver al CMAT en Cálculo cambia según el nivel de utilización observado durante MU?
\end{questionbox}

La probabilidad de cualquier visita durante Cálculo aumenta de manera marcada a través de los cuatro grupos de MU:

\begin{table}[H]
\centering
\caption{Probabilidad de utilizar el CMAT en Cálculo I según utilización previa durante MU.}
\label{tab:ppa_v8_persistence}
\begin{tabular}{@{}lrrrr@{}}
\toprule
Visitas MU & $N$ & Usuarios en Cálculo & Probabilidad & IC95\% Wilson \\
\midrule
0 & 2,470 & 426 & 17.25\% & [15.81\%, 18.79\%] \\
1--2 & 472 & 205 & 43.43\% & [39.03\%, 47.94\%] \\
3 & 103 & 56 & 54.37\% & [44.77\%, 63.66\%] \\
4+ & 196 & 131 & 66.84\% & [59.98\%, 73.05\%] \\
\bottomrule
\end{tabular}
\end{table}

La asociación global $4\times2$ es fuerte: $\chi^2(3)=392.47$, $p=9.47\times10^{-85}$, con $V$ de Cramér $=0.348$. La menor frecuencia esperada es 26.0, por lo que la aproximación asintótica de Pearson no presenta un problema de celdas escasas en esta tabla.

\begin{figure}[H]
\centering
\includegraphics[width=0.78\textwidth]{figures/15_ppa_persistence_by_mu_group}
\caption{Persistencia de uso del CMAT durante Cálculo I según número de visitas durante MU en la cohorte primaria de progresión. Los intervalos son Wilson al 95\%.}
\label{fig:ppa_v8_persistence}
\end{figure}

La comparación preespecificada más cercana a la pregunta ``cumplió exactamente el umbral'' versus ``continuó después del umbral'' es $V_{MU}=3$ frente a $V_{MU}\geq4$. La probabilidad observada de uso posterior es 54.37\% para exactamente tres visitas y 66.84\% para cuatro o más. La diferencia es
\[
RD=0.1247,
\qquad IC_{95\%}^{Newcombe}=[0.0092,0.2390].
\]
El riesgo relativo es
\[
RR=1.229,
\qquad IC_{95\%}=[1.004,1.505],
\]
y la razón de momios es $OR=1.691$ (IC95\% $[1.038,2.757]$). Pearson da $p=0.0343$ y Fisher bilateral $p=0.0440$.

\begin{answerbox}
Dentro de esta cohorte, quienes continúan más allá de tres visitas durante MU tienen una probabilidad posterior de volver al CMAT aproximadamente 12.5 puntos porcentuales mayor que quienes se detienen exactamente en tres. El resultado es compatible con una diferencia entre utilización limitada al umbral y utilización más persistente, pero no identifica por qué existe esa diferencia.
\end{answerbox}

Un modelo logístico con errores estándar agrupados por salón de MU conserva el gradiente. Tomando $V_{MU}=0$ como referencia y ajustando además por $Z_{MU}$, licenciatura oficial contemporánea al intento real de MU y periodo, los odds ratios estimados para utilizar CMAT posteriormente en Cálculo son:
\[
OR_{1-2}=3.53,
\qquad
OR_{3}=4.84,
\qquad
OR_{4+}=9.82.
\]
Los tres coeficientes son estadísticamente distintos de cero ($p<10^{-12}$ para 3 y 4+, y $p<10^{-28}$ para 1--2). Estos odds ratios no son riesgos relativos y no deben leerse como multiplicadores causales del PPA; muestran que la información contenida en el patrón previo de uso sigue siendo muy predictiva después de condicionar por desempeño relativo previo, carrera y periodo.

\subsubsection{¿Existe un cambio de comportamiento justo después de tres visitas?}

Como análisis descriptivo de forma funcional, se estima un logit por tramos:
\[
\operatorname{logit}P(C_i=1)
=\alpha+\beta_1\min(V_{MU},3)+\beta_2\max(V_{MU}-3,0),
\]
con el conteo truncado a 12 para que una cola muy escasa no domine la pendiente. Este modelo \textbf{no es una regresión discontinua}: el estudiante controla su número de visitas y no existe asignación exógena alrededor de $V=3$.

Hasta el umbral, el odds ratio por una visita adicional es $1.994$ (IC95\% $[1.807,2.200]$, $p=5.93\times10^{-43}$). Después del umbral, el odds ratio por visita adicional es $1.063$ (IC95\% $[0.982,1.150]$, $p=0.131$). La combinación de ambos resultados sugiere una relación muy pronunciada entre pasar de no uso/uso mínimo hacia el umbral y la persistencia futura, seguida por un aplanamiento de la pendiente. El contraste específico 3 versus 4+ todavía detecta una diferencia, pero los datos no respaldan una interpretación de incremento constante por cada visita posterior.

\subsection{Sensibilidad pedagógica: visitas específicamente etiquetadas como MU y Cálculo}

El PPA se relaciona con cualquier registro del CMAT durante el semestre, sin importar qué materia declaró el estudiante. Esa es la exposición correcta para estudiar el entorno del incentivo. Sin embargo, el Google Form permite una segunda pregunta más pedagógica: ¿buscar apoyo específicamente para MU predice buscar apoyo específicamente para Cálculo?

En la cohorte primaria, 660 estudiantes registraron al menos una asesoría etiquetada específicamente como MU. De ellos, 276 (41.82\%) registraron posteriormente al menos una asesoría etiquetada como Cálculo. Entre los 2,581 que no tuvieron asesoría etiquetada como MU, 365 (14.14\%) registraron apoyo específico de Cálculo. La diferencia es 27.68 puntos porcentuales (IC95\% $[23.72,31.69]$), con
\[
RR=2.957\;[2.594,3.371],
\qquad
OR=4.364\;[3.608,5.278].
\]
Pearson obtiene $\chi^2(1)=253.75$, $p=3.96\times10^{-57}$, con frecuencia esperada mínima 130.5; Fisher confirma el resultado ($p=3.33\times10^{-50}$).

Esta sensibilidad reduce la ambigüedad de que un estudiante haya visitado el CMAT por otra materia durante el semestre. El patrón de persistencia se mantiene cuando la exposición y el outcome se restringen a las materias específicas.

\subsection{¿Los patrones previos de uso se relacionan con desempeño posterior en Cálculo?}

La pregunta académica más defendible temporalmente no es si las visitas de Cálculo ``explican'' la nota de Cálculo, porque ambas ocurren durante el mismo curso y la dificultad experimentada puede aumentar simultáneamente las visitas. En cambio, el uso durante MU ocurre antes del outcome de Cálculo. Por ello se estima
\[
Z_{Calc,i}
=\beta_0+\beta_1Z_{MU,i}+\beta_2G^{MU}_i+\gamma L_{Calc,i}+\lambda_t+\varepsilon_i,
\]
con errores estándar agrupados por salón de Cálculo. No se incluye $V_{Calc}$ en la especificación principal.

Sin ajuste de licenciatura/periodo, pero controlando $Z_{MU}$, los coeficientes asociados a patrones previos de uso son $0.154Z$ para 1--2, $0.207Z$ para exactamente 3 y $0.200Z$ para 4+, todos respecto de cero visitas. Al añadir licenciatura contemporánea a Cálculo y periodo, las estimaciones son:
\[
\widehat\beta_{1-2}=0.122\;[0.051,0.194],\quad p=0.00081,
\]
\[
\widehat\beta_{3}=0.153\;[0.037,0.269],\quad p=0.00958,
\]
\[
\widehat\beta_{4+}=0.172\;[0.070,0.274],\quad p=0.000915.
\]
El coeficiente de desempeño previo es $\widehat\beta_{Z_{MU}}=0.578$ (IC95\% $[0.495,0.661]$).

\begin{answerbox}
El patrón de utilización observado durante MU contiene información sobre desempeño relativo posterior. Incluso entre estudiantes con $Z_{MU}$ similar, los grupos que habían utilizado el CMAT durante MU muestran, en promedio, un $Z_{Calc}$ posterior entre aproximadamente 0.12 y 0.17 desviaciones estándar mayor que el grupo de cero visitas después de ajustar por licenciatura y periodo. Este es un resultado longitudinal temporalmente ordenado, pero sigue siendo observacional: organización, motivación, preparación no medida y otras diferencias estables pueden explicar parte o toda la asociación.
\end{answerbox}

Los ocho perfiles muestran además diferencias descriptivas en $\Delta Z$. La Figura~\ref{fig:ppa_v8_delta} se conserva como exploración y no como prueba causal, porque la utilización durante Cálculo es contemporánea al segundo outcome.

\begin{figure}[H]
\centering
\includegraphics[width=0.84\textwidth]{figures/16_ppa_academic_trajectory_profiles}
\caption{Cambio medio en posición relativa $\Delta Z=Z_{Calc}-Z_{MU}$ por perfil observable de uso. La figura es descriptiva; el uso durante Cálculo puede ser respuesta a la dificultad experimentada en el propio curso.}
\label{fig:ppa_v8_delta}
\end{figure}

\subsection{La licenciatura sí forma parte de la heterogeneidad observada}

La carrera importa por dos razones distintas. Primero, diferentes programas tienen distinta exposición curricular a matemáticas y pueden generar diferente demanda de apoyo. Segundo, el archivo académico puede replicar materias aprobadas cuando un estudiante agrega o cambia de licenciatura, de modo que el programa debe vincularse al intento real y no a una revalidación posterior.

En la cohorte primaria, solo 11 de 3,241 pares (0.34\%) muestran un código oficial de licenciatura distinto entre el intento real de MU y el intento real posterior de Cálculo. No obstante, las tasas de uso y persistencia varían entre programas. Un logit ajustado por grupo de uso en MU, $Z_{MU}$ y periodo produce una prueba conjunta de los términos de licenciatura de
\[
W=97.43,\qquad df=24,\qquad p=8.23\times10^{-11},
\]
con errores estándar agrupados por salón de MU. Por tanto, la licenciatura contiene información adicional sobre la probabilidad de utilizar el CMAT durante Cálculo, incluso después de controlar el patrón previo de uso y el desempeño relativo de MU.

También existen diferencias en trayectoria académica relativa por licenciatura. Para $\Delta Z=Z_{Calc}-Z_{MU}$, Brown--Forsythe rechaza homogeneidad de varianzas entre las 24 licenciaturas con al menos 30 estudiantes ($p=2.43\times10^{-5}$). Por esa razón se utiliza ANOVA de Welch, que obtiene
\[
F_W(23,782.71)=6.246,\qquad p=1.39\times10^{-17}.
\]
Este resultado no significa que la licenciatura ``cause'' cambios de desempeño. Indica que el desplazamiento relativo entre MU y Cálculo no es homogéneo entre programas y justifica mantener la carrera como fuente de heterogeneidad observable en los modelos.

La licenciatura declarada en Google Forms se conserva como fuente paralela para estudiar la identidad académica reportada al momento de la asesoría. En esta cohorte aparece únicamente entre usuarios del CMAT (23.79\% durante MU y 25.24\% durante Cálculo, exactamente porque un formulario existe solo cuando hay visita). Por esta razón es inapropiado imputarla como si fuese un covariable basal para los no usuarios. El pipeline genera un cruce descriptivo entre código oficial y nombre declarado para futuras tareas de depuración/crosswalk.

\subsection{Tesis, hipótesis rival y síntesis de la nueva arquitectura}

\paragraph{Tesis compatible con los datos.}
El incentivo PPA1 puede reducir la barrera de entrada al CMAT durante el primer semestre. Una vez que el estudiante conoce el servicio, una parte de los usuarios continúa recurriendo al centro cuando llega a Cálculo I, aun cuando el mismo incentivo PPA1 ya no sea el motivo institucional relevante. La fuerte asociación entre uso previo y uso posterior, junto con la persistencia de apoyo específico MU $\rightarrow$ Cálculo, es compatible con un mecanismo de adopción y familiarización.

\paragraph{Hipótesis rival.}
La persistencia puede existir sin que el PPA haya creado el hábito. Estudiantes con una propensión estable a pedir ayuda, mayor organización, distinta preparación previa, diferentes exigencias de licenciatura o necesidades académicas recurrentes podrían utilizar el CMAT en ambos semestres. El incremento de persistencia desde 0 hasta 4+ visitas puede reflejar esa heterogeneidad individual. Los datos actuales no observan diagnóstico inicial, promedio de preparatoria ni motivación individual, por lo que no separan estos mecanismos.

\paragraph{Síntesis.}
La evidencia más útil no es una afirmación psicológica del tipo ``los de tres visitas solo querían puntos''. Los datos permiten una formulación operacional más precisa:
\[
\boxed{
\begin{aligned}
V_{MU}=3,\;V_{Calc}=0 &\quad\Rightarrow\quad \text{patrón compatible con uso limitado al umbral},\\
V_{MU}\geq4,\;V_{Calc}>0 &\quad\Rightarrow\quad \text{uso persistente más allá del contexto incentivado},\\
V_{MU}=0,\;V_{Calc}>0 &\quad\Rightarrow\quad \text{adopción tardía}.
\end{aligned}}
\]

Tres hallazgos deben conservarse para la futura reorganización del paper. Primero, la probabilidad de volver al CMAT aumenta de 17.25\% entre no usuarios de MU a 66.84\% entre usuarios de 4+. Segundo, exactamente 3 frente a 4+ presenta una diferencia posterior de 12.5 puntos porcentuales, compatible con una distinción conductual alrededor del umbral, aunque no causal. Tercero, el patrón de uso previo sigue asociado con desempeño relativo posterior en Cálculo aun después de condicionar por $Z_{MU}$, licenciatura y periodo.

\begin{cautionbox}[title=Alcance inferencial del capítulo]
Ninguno de estos resultados identifica el efecto causal de PPA1 ni el efecto causal de las asesorías. El número de visitas es elegido por el estudiante; el umbral $V=3$ no es una variable de asignación exógena y no justifica un diseño de regresión discontinua. La selección de la cohorte también condiciona en haber aprobado MU y haber llegado a Cálculo, de modo que las conclusiones longitudinales se refieren a progresores, no a todos los estudiantes que iniciaron Matemáticas Universitarias.
\end{cautionbox}


\clearpage
\appendix
\section{Resultados estadísticos complementarios}

\subsection{Contraste PPA alcanzado: $V\geq3$ frente a $V<3$}

El contraste alternativo que pregunta por alcanzar el umbral --en vez de continuar más allá de él-- produce una diferencia media no ajustada de $0.362Z$ (\CI\ $[0.290,0.429]$), probabilidad de superioridad $0.617$ y Cliff $\delta=0.234$. Con profesor-periodo y carrera, el coeficiente es $0.382$ (\CI\ $[0.287,0.476]$). La similitud con el contraste $V>3$ confirma que ambos indicadores capturan, en gran medida, la separación entre usuarios recurrentes y la mayoría de estudiantes con cero o pocas visitas.

\subsection{Tabla de sensibilidad del outcome continuo}

\begin{table}[H]
\centering
\caption{Sensibilidad del contraste $V>3$ frente a $V\leq3$.}
\begin{tabular}{@{}lrrr@{}}
\toprule
\textbf{Outcome} & \textbf{Diferencia media} & \textbf{\CI} & \textbf{Cliff $\delta$} \\
\midrule
Imputación primaria adversa & 0.359 & $[0.277,0.439]$ & 0.233 \\
Imputación uniforme bajo 7.5 & 0.343 & $[0.257,0.424]$ & 0.227 \\
Solo calificaciones numéricas & 0.211 & $[0.137,0.281]$ & 0.075 \\
\bottomrule
\end{tabular}
\end{table}

\subsection{Regularidad temporal entre quienes alcanzan PPA}

\begin{table}[H]
\centering
\caption{Meses activos entre estudiantes con $V\geq3$.}
\begin{tabular}{@{}rrrr@{}}
\toprule
\textbf{Meses activos} & \textbf{$N$} & \textbf{Visitas medias} & \textbf{Semanas efectivas medias} \\
\midrule
1 & 43 & 3.53 & 2.08 \\
2 & 156 & 3.85 & 2.89 \\
3 & 154 & 4.77 & 3.74 \\
4 & 119 & 9.07 & 6.43 \\
\bottomrule
\end{tabular}
\end{table}

\section{Figura temporal completa}

\begin{figure}[H]
\centering
\includegraphics[width=\textwidth]{figures/05_daily_cmat_timeline.png}
\caption{Serie diaria completa de actividad del CMAT. La línea suavizada de siete días muestra periodos recurrentes de aumento y descenso de demanda.}
\end{figure}


\begin{thebibliography}{99}
\bibitem[Mullen et~al.(2024)Mullen, Howard, and Cronin]{mullen2024}
Mullen, C., Howard, E., and Cronin, A. (2024).
\newblock A scoping literature review of the impact and evaluation of mathematics and statistics support in higher education.
\newblock \emph{Educational Studies in Mathematics}, 117, 1--22. doi:10.1007/s10649-024-10332-6.

\bibitem[Lawson et~al.(2020)Lawson, Grove, and Croft]{lawson2020}
Lawson, D., Grove, M., and Croft, T. (2020).
\newblock The evolution of mathematics support: a literature review.
\newblock \emph{International Journal of Mathematical Education in Science and Technology}, 51(8), 1224--1254.

\bibitem[Rickard and Mills(2018)]{rickard2018}
Rickard, B., and Mills, M. (2018).
\newblock The effect of attending tutoring on course grades in Calculus I.
\newblock \emph{International Journal of Mathematical Education in Science and Technology}, 49(3), 341--354.

\bibitem[Mac an Bhaird et~al.(2009)Mac an Bhaird, Morgan, and O'Shea]{macanbhaird2009}
Mac an Bhaird, C., Morgan, T., and O'Shea, A. (2009).
\newblock The impact of the mathematics support centre on the grades of first year students at the National University of Ireland Maynooth.
\newblock \emph{Teaching Mathematics and its Applications}, 28(3), 117--122.

\bibitem[Jacob and N\'{i} Fhloinn(2019)]{jacob2019}
Jacob, M., and N\'{i} Fhloinn, E. (2019).
\newblock A quantitative, longitudinal analysis of the impact of mathematics support in an Irish university.
\newblock \emph{Teaching Mathematics and its Applications}, 38(4), 216--229.

\bibitem[B\"uchele and Sch\"urmann(2024)]{buchele2024}
B\"uchele, S., and Sch\"urmann, M. (2024).
\newblock Causal evidence of the effect of a course specific mathematical learning and support center on economics students' performance and affective variables.
\newblock \emph{Studies in Higher Education}, 49(9), 1569--1587.

\bibitem[Paloyo et~al.(2016)Paloyo, Rogan, and Siminski]{paloyo2016}
Paloyo, A. R., Rogan, S., and Siminski, P. (2016).
\newblock The effect of supplemental instruction on academic performance: An encouragement design experiment.
\newblock \emph{Economics of Education Review}, 55, 57--69.

\bibitem[Pugatch and Wilson(2018)]{pugatch2018}
Pugatch, T., and Wilson, N. (2018).
\newblock Nudging study habits: A field experiment on peer tutoring in higher education.
\newblock \emph{Economics of Education Review}, 62, 151--161.

\bibitem[Fong et~al.(2023)Fong, Gonzales, Hill-Troglin Cox, and Shinn]{fong2023}
Fong, C. J., Gonzales, C., Hill-Troglin Cox, C., and Shinn, H. B. (2023).
\newblock Academic help-seeking and achievement of postsecondary students: A meta-analytic investigation.
\newblock \emph{Journal of Educational Psychology}, 115(1), 1--21.

\bibitem[Rosenbaum and Rubin(1983)]{rosenbaum1983}
Rosenbaum, P. R., and Rubin, D. B. (1983).
\newblock The central role of the propensity score in observational studies for causal effects.
\newblock \emph{Biometrika}, 70(1), 41--55.

\bibitem[Austin(2011)]{austin2011}
Austin, P. C. (2011).
\newblock An Introduction to Propensity Score Methods for Reducing the Effects of Confounding in Observational Studies.
\newblock \emph{Multivariate Behavioral Research}, 46(3), 399--424.

\bibitem[Hern\'{a}n and Robins(2020)]{hernan2020}
Hern\'{a}n, M. A., and Robins, J. M. (2020).
\newblock \emph{Causal Inference: What If}.
\newblock Chapman \& Hall/CRC.
\end{thebibliography}

\end{document}
